%%%%%%%%%%%%%%%%%%%% author.tex %%%%%%%%%%%%%%%%%%%%%%%%%%%%%%%%%
%
% sample root file for your "contribution" to a contributed volume
%
% Use this file as a template for your own input.
%
%%%%%%%%%%%%%%%% Springer Nature %%%%%%%%%%%%%%%%%%%%%%%%%%%%%%%%


% RECOMMENDED %%%%%%%%%%%%%%%%%%%%%%%%%%%%%%%\documentclass[runningheads]{llncs}
\documentclass[conference,compsoc]{IEEEtran}

\ifCLASSOPTIONcompsoc
  % IEEE Computer Society needs nocompress option
  % requires cite.sty v4.0 or later (November 2003)
  \usepackage[nocompress]{cite}
\else
  % normal IEEE
  \usepackage{cite}
\fi


\usepackage[T1]{fontenc}
\usepackage{graphicx}
\usepackage{booktabs}
\usepackage{microtype}
\usepackage{pgfplots}
\usepackage{pgfplotstable}
\usepackage{float}
\usepackage{caption}
\usepackage{titlesec}
\usepackage{multirow}
% \setcounter{secnumdepth}{3}
% \titlespacing*{\section}
% {3pt}{1.2ex plus .2ex}{0.8ex}

% \titlespacing*{\subsection}
% {3pt}{1ex}{0.5ex}

% \usepgfplotslibrary{statistics}
% \pgfplotsset{compat=1.18}
\usepackage{pgfplots}
\usepgfplotslibrary{statistics}
\pgfplotsset{compat=1.18}

\def\BibTeX{{\rm B\kern-.05em{\sc i\kern-.025em b}\kern-.08em
    T\kern-.1667em\lower.7ex\hbox{E}\kern-.125emX}}

\makeatletter
\newcommand{\linebreakand}{%
  \end{@IEEEauthorhalign}
  \hfill\mbox{}\par
  \mbox{}\hfill\begin{@IEEEauthorhalign}
}
\makeatother
\begin{document}
% \motto{This is my motto.}
\title{Leveraging Historical CWE Exploitability and Bi-Encoder-based CWE Classification for Early Vulnerability Prioritization}
CVE -> CWE

\author{

\IEEEauthorblockN{Person A}
\IEEEauthorblockA{
    \textit{Some University} \\
    email@email.com
}

\and
\IEEEauthorblockN{Person B}
\IEEEauthorblockA{
    \textit{Some University} \\
    email@email.com
}

\and
\IEEEauthorblockN{Person C}
\IEEEauthorblockA{
    \textit{Some University} \\
    email@email.com
}

\and
\IEEEauthorblockN{Person D}
\IEEEauthorblockA{
    \textit{Some University} \\
    email@email.com
}
}
% \author{
% Kobra Khanmohammadi\inst{1}\orcidID{0009-0004-1414-2111}  \and
% Alexander Da Re\inst{1} \and Benjamin Whelan\inst{1} \and
% Raphaël Khoury\inst{2}\orcidID{0000-0002-7625-3384} 
% }

% \authorrunning{K. Khanmohammadi et al.}

% \institute{
% Sheridan College, Ontario, Canada\\
% \email{kobra.khanmohammadi@sheridancollege.ca, dareal@sheridancollege.ca, whelaben@sheridancollege.ca}
% \and
% Université du Québec en Outaouais (UQO), Canada\\
% \email{raphael.khoury@uqo.ca}}
%
% Use the package "url.sty" to avoid
% problems with special character[s
% used in your e-mail or web address
%
\maketitle


% \begin{abstract}
% Prioritizing newly disclosed software vulnerabilities is challenging because limited information is available immediately after disclosure. While the Common Vulnerability Scoring System (CVSS) measures vulnerability severity, it does not estimate the likelihood of real-world exploitation. In this paper, we investigate whether the historical exploitability of Common Weakness Enumeration (CWE) categories provides a reliable empirical basis for supporting the early prioritization of newly disclosed Common Vulnerabilities and Exposures (CVEs). We first conduct a large-scale empirical study of 192,788 CVEs, 768 CWE categories, and exploitation evidence from the CISA Known Exploited Vulnerabilities (KEV), VulnCheck, and GrayNoise datasets. Our analysis characterizes the evolution of CWE mappings and demonstrates that historical exploitability patterns of CWE categories exhibit statistically significant temporal stability, and identifies software weaknesses consistently associated with exploitation in the wild. Motivated by these findings, we propose a hierarchy-aware Large Language Model (LLM) framework that automatically maps newly disclosed CVEs to CWE categories, enabling practitioners to leverage historical weakness-level exploitability information immediately after vulnerability disclosure. The framework leverages the hierarchical structure of the CWE taxonomy and contextual information extracted from CVE reference links to improve classification accuracy. Experimental results show that taxonomy-aware prompting, domain-specialized LLMs, and contextual enrichment improve CWE classification, these findings establish the empirical and technical foundations for weakness-informed vulnerability prioritization before CVE-specific exploitability evidence becomes available.
% \keywords{Common Vulnerabilities and Exposures (CVE) \and Common Weakness Enumeration (CWE) \and Large Language Models \and Vulnerability Prioritization \and Exploitability Analysis}
% \end{abstract}

%%%%%%%%%%%%%%%%%%%%%%%
%%%%
\begin{abstract}The prioritization of newly discovered vulnerabilities can be challenging because crucial information is often missing in the period that initially follows disclosure. Metric such as the Common Vulnerability Scoring System (CVSS), even if available, does not strongly correlate with the actual likelihood of exploitation. In this paper, we investigate whether the historical exploitability of Common Weakness Enumeration (CWE) categories provides a reliable empirical basis for supporting the early prioritization of newly disclosed Common Vulnerabilities and Exposures (CVEs). We first conduct a large-scale empirical study of 192,788 CVEs, 768 CWE categories, and exploitation evidence from several sources.  Our analysis shows that CWE mappings are stable over time and can serve as a strong predictor of exploitability. This finding motivates, in turn, the creation of a fine-tuned bi-encoder model that automatically maps newly disclosed CVEs to CWE categories by first predicting the main category and then the child, enabling practitioners to leverage historical weakness-level exploitability information immediately after vulnerability disclosure. Experimental results show that this framework enhances CWE categorization and can be used to accurately forecast CVE exploitability.
% \keywords{Common Vulnerabilities and Exposures (CVE) \and Common Weakness Enumeration (CWE) \and Large Language Models \and Vulnerability Prioritization \and Exploitability Analysis} 
\end{abstract}
%%%%%%%%%%%%%%%%%%%%%%%%%%%%%%%%%%%%%%%%%%%%

\section{Introduction}

Modern organizations face an unprecedented number of publicly disclosed software vulnerabilities. Every month, the National Vulnerability Database (NVD) publishes thousands of new CVEs, far more than security teams can address at once. Deciding which vulnerabilities require immediate attention is therefore a critical problem, and it must often be solved within days of disclosure, before reliable exploitation signals exist.

Existing prioritization approaches fall short in this early window. CVSS severity scores are poor predictors of exploitation, and exploit-prediction methods operate at the level of individual CVEs and depend on evidence that accumulates only after disclosure. What is missing is a signal that is available \emph{at} disclosure and that generalizes across vulnerabilities.

We propose that such a signal can come from the underlying software weakness. Vulnerabilities that share a CWE often share exploitation characteristics, so the historical exploitation behavior of a CWE may indicate the risk of a new CVE that maps to it. Realizing this idea requires knowing the CWE of a new CVE, which is often unassigned at publication. We therefore study both whether weakness-level exploitability is a usable signal and whether the CWE can be inferred automatically from the CVE description. Concretely, this paper addresses three research questions:

\begin{itemize}
    \item[\textbf{RQ1:}] Are historical exploitability patterns at the CWE level stable enough to provide actionable predictive information?
    \item[\textbf{RQ2:}] Can newly disclosed CVEs be accurately mapped to their underlying CWE automatically?
    \item[\textbf{RQ3:}] Can weakness-level information improve exploitability prediction?
\end{itemize}

To answer these questions, this work makes the following contributions:
\begin{itemize}
    \item A large-scale empirical study of 192,788 CVEs and 768 CWE categories, characterizing how CWE assignments evolve and how weaknesses relate to real-world exploitation. The study shows that historical exploitability at the CWE level is temporally stable, with certain weakness categories consistently showing elevated exploitation rates.
    \item A hierarchy-aware bi-encoder framework that maps CVE descriptions to CWE categories by predicting the parent category first and then the child.
    \item A demonstration that automatically predicted CWEs enable weakness-informed exploitability forecasting immediately after disclosure.
\end{itemize}

% TODO: fill in headline results (e.g., classification accuracy, exploitability prediction performance).
Our results show that CWE-level exploitability is stable over time and that the proposed framework improves CWE classification and supports accurate exploitability forecasting.

% The remainder of this paper is organized as follows. Section~\ref{sec:motivation} motivates the approach, Section~\ref{sec:study} presents the empirical study, Section~\ref{sec:method} describes the classification framework, and Section~\ref{sec:eval} reports the evaluation.



\section{Background and Motivation}
\label{sec:motivation}


\subsection{Limits of Severity and CVE-Level Prioritization}
Although CVSS remains the de facto standard for measuring vulnerability severity, numerous studies have shown that CVSS scores alone are poor predictors of whether a vulnerability will be exploited in the wild \cite{jacobs2023epss,sherif2026kri,jiang2025survey,khoury2025mining}. Vulnerability management has consequently shifted toward exploitability-aware approaches, such as the Exploit Prediction Scoring System (EPSS), which estimate the likelihood of future exploitation \cite{jacobs2023epss}. These approaches operate at the individual CVE level and rely on evidence that is often missing or unstable right after disclosure, such as public exploit code, external references, and community discussion. % TODO: adjust this list to the features you actually mean.
This limits their usefulness precisely when prioritization decisions are most urgent.

\subsection{CWE as an Early Signal}
Every vulnerability originates from one or more underlying software weaknesses described by the CWE. Because vulnerabilities sharing a weakness often exhibit similar exploitation characteristics, the historical exploitation behavior of a CWE can serve as an early signal of risk for a new CVE, before the features required by CVE-level predictors become available. % TODO: state how CVEs with multiple CWEs are handled (e.g., primary CWE only, or multi-label).
Whether such weakness-level patterns remain stable over time, and therefore whether they can support early prioritization, is the question addressed by RQ1.

\subsection{Obstacles to Using CWE in Practice}
Two challenges stand in the way of using weakness-level information in practice.
First, \emph{timing and scale}: newly published CVEs are frequently released before an analyst assigns an official CWE, and manual assignment requires expert analysis that does not scale with the growing volume of disclosures.
Second, \emph{label ambiguity}: the CWE taxonomy is hierarchical \cite{Bojanova}, with weaknesses ranging from broad categories to specific variants, and experts often disagree on the most appropriate CWE for a given vulnerability. Any automated classifier must cope with this ambiguity.

\subsection{Embedding-Based CWE Classification}
Our approach is embedding-based CWE classification. A bi-encoder encodes CVE descriptions and CWE entries into a shared vector space, and a CVE is mapped to the CWE whose embedding is closest to its description. Such architectures have performed well on technical security text and on classifying vulnerabilities into CWE categories \cite{das2021v2w,alhaddad2025prompting}. This design addresses both obstacles above. To address timing and scale, the CWE is inferred directly from the CVE description, so no analyst input is needed at disclosure. To address label ambiguity, the model exploits the hierarchy of the taxonomy: it first predicts the parent category and then refines the prediction to a child, so that disagreements between closely related weaknesses do not obscure the broader category. The predicted CWE is then linked to its historical exploitability, making weakness-level risk information available at the moment of disclosure.

\section{Literature Review}

Vulnerability prioritization has traditionally relied on the CVSS, even though severity scores alone have repeatedly been shown to be poor indicators of real world exploitation. The EPSS of Jacobs et al.~addresses this limitation by estimating the probability that an individual CVE will be exploited~\cite{jacobs2023epss}, and subsequent work has extended this direction: Sherif et al.~combine threat, impact, and exposure into a Key Risk Indicator framework~\cite{sherif2026kri}, Khanmohammadi et al.~propose \textit{ExploitabilityBirthMark}, which predicts exploitability from information available at disclosure~\cite{khanmohammadi2025exploitabilitybirthmark}, and Jiang et al.~survey the broader shift toward dynamic prioritization methods that account for exploitability~\cite{jiang2025survey}. These methods, however, all reason about individual CVEs; even when the assigned CWE appears among their predictive features, none examines the exploitation history of the weakness categories themselves.

A second line of work moves from vulnerabilities to weaknesses, developing metrics that rank CWE categories by importance. Mell and Gueye score weaknesses along dimensions such as vulnerability frequency, severity, impact, and exploitability~\cite{mell2020suite}, and Galhardo et al.~refine these rankings by correcting biases in the MITRE CWE Top 25 methodology and balancing exploitability against impact~\cite{galhardo2021measurements,cwe_top25}. Such rankings show which weaknesses matter most in aggregate, but they are descriptive rather than predictive: they offer no guidance for a specific, newly disclosed CVE.

A third line of work automates the mapping of CVE descriptions to CWE categories. Das et al. \cite{das2021v2w} proposed V2W-BERT,
a hierarchical BERT-based classification framework,
reporting up to 97\% accuracy with random data partitioning and 94\% with temporal partitioning. More recently, Simonetto et al. \cite{simonetto2026knowing} leveraged the CWE hierarchy and fine-tuned language models to improve fine-grained weakness classification, reporting an approximately five-percentage-point improvement in F1-score over a supervised baseline. However, these approaches rely on supervised classification and may face challenges with infrequent CWE classes and evolving vulnerability patterns.

Haddad et al. \cite{haddad2023automated} approached CVE-to-CWE mapping as a semantic ranking problem,
evaluating Sentence-BERT and rankT5 models on a
manually annotated dataset of 4,012 CVE records.
Their results demonstrated improvements over traditional information retrieval and transformer-based baselines.
Alternatively, Liu et al. \cite{liu2023not} investigated ChatGPT-based vulnerability mapping using 13,513 CVEs, reporting exact-match accuracies of 50.88\% and 53.13\% under two prompting strategies. When considering the CWE hierarchy, parent-level agreement increased to 72.20\% and 76.13\%, respectively. These findings highlight the difficulty of predicting exact CWE identifiers, even when the broader weakness category is correctly identified.

Despite these advances, several challenges remain,
including distinguishing semantically related CWEs,
handling imbalanced weakness distributions, and
maintaining prediction accuracy on newly disclosed
vulnerabilities. Moreover, supervised classifiers
often require substantial labeled training data,
while generative LLMs introduce inference costs and
potentially invalid CWE predictions. To address these
challenges, our study investigates a lightweight
semantic matching approach based on all-MiniLM-L6-v2,
fine-tuned using Low-Rank Adaptation (LoRA). Unlike
generative approaches, our method selects predictions
from a predefined set of CWE candidates, avoiding
the generation of invalid identifiers. We further
evaluate the model using temporally separated CVE
records to examine its generalization to future
vulnerability disclosures. 

% Our work thus investigates whether parameter-efficient semantic representations can provide an effective alternative to larger generative models and conventional supervised classification frameworks.

Taken together, these three lines of work remain largely disconnected: exploitability prediction operates on individual CVEs, weakness rankings describe CWEs in aggregate, and automated mapping assigns labels without linking them to exploitation outcomes. To the best of our knowledge, no prior study has examined whether historical exploitability at the CWE level is stable enough to serve as an early prioritization signal, nor shown how CWE classification that exploits the category hierarchy can put that signal to use in the first days after disclosure. The present work addresses both questions.




\section{Overview of the Proposed Framework}
\label{sec:overview}

The objective of this study is to investigate whether
historical exploitation patterns associated with
software weaknesses can support the early prioritization
of newly disclosed vulnerabilities. Unlike conventional
exploitability prediction approaches that primarily
rely on characteristics and threat intelligence
associated with individual CVEs, our approach leverages
the historical exploitation behavior of their underlying
CWE categories. The central hypothesis is that
vulnerabilities sharing a common weakness exhibit
partially consistent exploitation patterns, allowing
historical CWE-level information to provide an early
indication of exploitation risk.

To investigate this hypothesis, we propose a framework
consisting of three interconnected phases, as illustrated
in Figure~\ref{fig:framework_overview}.


\begin{figure*}[!t]
    \centering
    \includegraphics[width=\textwidth]{Figures/overview1.jpg}
    \caption{Overview of the proposed CWE-informed vulnerability
    prioritization framework.}
    \label{fig:framework_overview}
\end{figure*}


\textbf{Phase 1: Historical CWE Exploitability Analysis.}
We first conduct a large-scale empirical analysis of
CVE--CWE mappings and exploitation evidence collected
from NVD, CISA KEV, VulnCheck, and GrayNoise.
This analysis examines the availability and stability
of CWE assignments, the temporal consistency of
weakness-level exploitation patterns, and the
association between individual CWE categories and
observed exploitation in the wild. Based on these
findings, we construct historical exploitability
profiles that characterize the exploitation behavior
of software weaknesses.

\textbf{Phase 2: Automated CVE-to-CWE Classification.}
Because newly disclosed vulnerabilities may not have
an assigned CWE, the second phase introduces a
hierarchy-aware bi-encoder model to infer the underlying
weakness directly from the CVE description.
The model is based on all-MiniLM-L6-v2 and fine-tuned
using Low-Rank Adaptation (LoRA). It represents CVE
descriptions and CWE definitions in a shared embedding
space and uses semantic similarity to identify
candidate weaknesses. By exploiting the CWE taxonomy,
the model first identifies relevant high-level
categories and subsequently refines the prediction
toward more specific CWE identifiers. This process
enables automated weakness identification without
requiring an official CWE assignment.

\textbf{Phase 3: CWE-Informed Exploitability
Prioritization.}
In the final phase, the predicted CWE is linked to
its historical exploitability profile to estimate
the exploitation risk of the newly disclosed CVE.
Historical exploitation rates are estimated from
previously disclosed vulnerabilities using a
smoothed probability calculation to reduce the
influence of sparsely represented weaknesses.
These probabilities are used to rank newly disclosed
CVEs according to their estimated exploitation risk.
The framework is evaluated using temporally separated
data to examine whether historical weakness-level
information can support prioritization of future
vulnerabilities.

Together, these phases connect historical
weakness-level exploitation knowledge with automated
vulnerability classification and early exploitability
assessment. The empirical analysis establishes the
motivation for using CWE history, the bi-encoder
enables the identification of weaknesses for newly
disclosed CVEs, and the prioritization component
translates this information into actionable
vulnerability rankings.

\section{Empirical Study}
\label{sec:empirical_study}
\subsection{Requirements \& Approaches (Pooria)}
In order to create a system capable of assigning CWE-IDs to newly released CVEs it would have to qualify for the following requirements:
\begin{itemize}
    \item \textbf{Adaptable to new CWEs:} New CWE categories are rare, but they do occur and having a system that is capable of distinguishing new categories without retraining is crucial.
    \item \textbf{Independent of syntax and language:} CVEs often include product names, versions, and irrelevant information that could throw off the system. Therefore, the system should not rely on syntax, hard-coded rules, and keywords during inference.
    \item \textbf{No memorization:} The system should not see the test data during training and, ideally, should not memorize the CWE-IDs to CVE descriptions.
\end{itemize}
There are various ways to implement a mapping from CVE descriptions to CWE categories:
\begin{itemize}
    \item \textbf{LLM based:} LLMs are really good at domain shift applications, and our end goal matches that pattern very well since it's only transforming a sentence into another form. Several approaches [CITATION NEEDED] have been used to fine-tune custom LLMs for this purpose, and frontier LLMs seem to be very well versed in this domain. Nonetheless, it is very hard to determine whether an opensource LLM model contains our dataset as part of it's training data. Hence, we have to rely on LLMs trained on internet data before 2022.
    
    \item \textbf{BERT based:} BERT classifiers offer an effective solution to this problem as they can one-shot the solution and are easily fine-tuned. Our baseline uses a 125M parameter bert model from \cite{mosievskiy2026roberta} which achieves 87.4\% accuracy across our test dataset, which includes CVEs from 2021-2025
    \item \textbf{Keyword based:} Keyword matching systems are very useful in simple categorization examples where specific approaches are mentioned, such as SQL Injection, XSS, CSRF, CRYPTO, HTTP/HTTPS. However, they fail when no keyword is found or give false positives when examples are too similar. Furthermore, they do not provide a measure of confidence.
    \item \textbf{Embedding based:} Bi-Encoders or embedding based approaches are an underutilized method in this space; they are used very often in Retrieval-Augmented Generation (RAG) applications, which is very similar to our CVE $\leftrightarrow$ CWE matching problem. Unlike BERT based approaches, embeddings can provide a relative distance from other matches, which gives a proxy to confidence in this sense and lets us evaluate how the model is performing as it is training.
\end{itemize}
\subsection{Methodology (Pooria)}

% Methodology figure
\begin{figure}[H]
\centering
\includegraphics[width=1\columnwidth]{Figures/methodology/methodology.png}
\caption{CVE $\rightarrow$ CWE Mapping Methodology}
\label{fig:methodology}
\end{figure}

In this study, we use an embedding based approach to map CVE to CWE descriptions shown in Figure \ref{fig:methodology}. This approach allows us to employ a pipeline similar to RAG, where we use cosine similarity to match descriptions together; however, \textit{we do not use} LLMs to generate the final match. In this paper, we finetune an \texttt{all-MiniLM-L6-v2} model on our dataset from 1999-2021 CVEs and achieve state-of-the-art performance.

\subsubsection{CWE Dataset: }

We gather all of our CWE data from \texttt{cwe.mitre.org} under the \texttt{CWE-699} view, \textit{Software Development}. 
We scrape everything from parents to children and visualize their relationship in Figure \ref{fig:cwe_dependency_graph}, which, most importantly, shows that some weaknesses are under multiple parents. Children with multiple parents pose a problem because we cannot assume that each child has a single parent in our multistage search when finding the right CWE for a given CWE description. In order to solve this issue, we use BFS to first choose the parent that the model perceives as most appropriate, then fetch the children associated with that parent. % in order to isolate for noise and reduce candidates.

\begin{figure}[H]
\centering
\includegraphics[width=1\columnwidth]{Figures/methodology/cwe_dependency_graph.png}
\caption{CWE-699 Dependency Graph}
\label{fig:cwe_dependency_graph}
\end{figure}

\subsubsection{CVE Dataset: }


We use a training dataset consisting of CVEs from 1999-2021 and validation dataset of  CVEs from 2023-2025. The range for our validation dataset is very large which lets us evaluate the fine-tuned model extensively. Our training dataset is based off of \texttt{zefang-liu}'s dataset \cite{liu2024cvecwedataset}, to which we add CWE descriptions and perform further post-processing to remove as much noise as possible from the descriptions. 


\subsubsection{Training methods}

We first considered prompting a large language model to output the CWE for a CVE description. Prompting needs no training data and can be given the CWE hierarchy in the prompt, but it has two problems in our setting. First, we cannot verify whether a model saw a given CVE or its CWE label during pretraining, so its evaluation is unreliable. Second, a generated label gives no calibrated distance, so we cannot rank the alternatives or judge how confident the model is. We therefore chose an encoder that embeds the CVE description and the CWE text into one vector space and returns a similarity score for every candidate. The score gives a natural ranking and a confidence proxy, and the encoder can be finetuned on our own data, which removes the contamination concern.

We then chose how to finetune. Full finetuning updates every weight of the encoder and is prone to overfitting here, because the descriptions repeat many phrases across CVEs and the label space is small. Low Rank Adaptation (LoRA) freezes the pretrained weights and learns a small number of low rank updates, which reduces overfitting, finishes in a few hours on a single GPU, and lets us run both recipes under identical conditions. Because the adapters are merged back into the base weights, inference costs the same as the untuned encoder.

With the adaptation method fixed, we built the training data, following the pipeline in Figure~\ref{fig:methodology}. We split the CVEs from 1999 to 2021 by CVE identifier, so a CVE with several CWE labels never appears in two splits, and we stratified the split on the primary CWE so that rare weaknesses appear in validation and test. We rendered every CWE node as a short text that joins its title, its summary, and the titles of its immediate children. We trained with MultipleNegativesRankingLoss, where each (anchor, positive) pair is contrasted against the other positives in the batch, and we used the NO\_DUPLICATES sampler so that no negative equals the positive. We then mined difficult negatives with the untuned encoder: for each CWE we selected the nearest nodes that are neither related nor siblings, that is, text that looks similar but belongs to another category, and the nearest categories that are not related. Finally we formed two recipes from these pairs. The children only recipe uses plain (CVE, own CWE) pairs with the batch negatives only. The hierarchy full recipe adds three sets under the same loss: CVE to child CWE with one sibling and one similar CWE from another category as negatives, CVE to parent category with the nearest other category as negative, and a CWE to sibling geometry set that pulls siblings together and pushes similar nodes from other categories apart. Both recipes share the split, the base model, the LoRA initialization, the optimizer, and the seed, so the comparison isolates the effect of the training pairs.


\subsection{Model}

\subsubsection{Baseline}

We evaluate against two baselines. The first is the untuned \texttt{all-MiniLM-L6-v2} encoder used with the same retrieval procedure, which measures the contribution of finetuning. The second is the published flat classifier of Mosievskiy et al.~\cite{mosievskiy2026roberta}, a 125M parameter RoBERTa model that predicts CWE labels directly and reports 87.4\% top 1 accuracy on its own benchmark. Since it returns only a flat ranking of children, we map each predicted child to its parent to obtain a category ranking.

\subsubsection{Model fine-tuning}

We finetune \texttt{all-MiniLM-L6-v2} with LoRA (rank 64, alpha 128, dropout 0, target modules \texttt{value}, \texttt{key}, \texttt{dense}, \texttt{query}). Training uses 3 epochs, learning rate $1\times10^{-4}$, batch size 64, a linear schedule with 5\% warmup, mixed precision, and seed 42. We choose the model on the portion of the 1999 to 2021 CVEs that we set aside, and we merge the adapters into the base weights before evaluation, so inference is one forward pass of the encoder.

\subsubsection{Results}

Table~\ref{tab:recall} reports Recall@$k$ on the CVEs disclosed from 2023 to 2025, which the models did not see during training, and Figure~\ref{fig:recall} shows the corresponding curves.

Finetuning improves every target over the untuned encoder. The untuned baseline recovers only a small share of the correct categories and children, while both recipes improve on it by a wide margin, which shows that the retrieval signal can be learned from our data.

The flat child target gives the highest recall of the child targets. Here the model searches the whole taxonomy directly rather than only the children of the top categories, so it is the most favorable and the most direct test of child prediction. Even on this target, recipe A (hier full), which is trained with the structure of the taxonomy, finds more correct children than recipe B (children only), which is trained on flat pairs. Training with the hierarchy therefore improves child prediction in general and not only when the search is restricted to a few categories. Recipe A also finds more correct children than the SOTA classifier.

The category stage explains this difference. Recipe A learns to place a CVE near its parent category, while recipe B does not, so recipe B loses accuracy both when the search is restricted to the top categories and, to a smaller degree, in the flat setting.

On the category target the SOTA classifier is slightly better for the first few guesses, but recipe A catches up as more guesses are allowed and the two are close by $k=5$. Overall, using the hierarchy during training gives the strongest child predictions, which is the quantity that matters for the exploitability step.

\begin{table}[H]
\centering

\caption{CVE$\rightarrow$CWE Recall@$k$ for CVEs disclosed from 2023 to 2025, which the models did not see during training. Best value in each column in bold. The published classifier has no category stage, so it cannot produce hierarchical child recall.}
\label{tab:recall}
\begin{tabular}{llccc}
\toprule
 &  & @1 & @3 & @5 \\
\midrule
\multirow[c]{4}{*}{Category} & Baseline (MiniLM) & 0.236 & 0.409 & 0.510 \\
 & A: hier full & 0.750 & 0.856 & \textbf{0.889} \\
 & B: children only & 0.584 & 0.776 & 0.841 \\
 & SOTA (cwe-predictor) & \textbf{0.802} & \textbf{0.862} & 0.869 \\
\cline{1-5}
\multirow[c]{3}{*}{Child (hier.)} & Baseline (MiniLM) & 0.124 & 0.175 & 0.216 \\
 & A: hier full & \textbf{0.695} & \textbf{0.735} & \textbf{0.756} \\
 & B: children only & 0.506 & 0.557 & 0.584 \\
\cline{1-5}
\multirow[c]{4}{*}{Child (flat)} & Baseline (MiniLM) & 0.337 & 0.432 & 0.479 \\
 & A: hier full & \textbf{0.753} & \textbf{0.828} & \textbf{0.853} \\
 & B: children only & 0.715 & 0.795 & 0.823 \\
 & SOTA (cwe-predictor) & 0.655 & 0.753 & 0.779 \\
\cline{1-5}
\bottomrule
\end{tabular}
\end{table}

\begin{figure}[H]
\centering
\includegraphics[width=1\columnwidth]{Figures/model/overall.png}
\caption{CVE$\rightarrow$CWE Recall@$k$ for the untuned baseline, the two finetuned recipes, and the published classifier. One panel per target.}
\label{fig:recall}
\end{figure}




\section{Empirical Study}
\subsection{Dataset and Methodology}
\label{sec:dataset_methodology}

This empirical study combines vulnerability information from multiple publicly available sources to investigate the evolution of CWE mappings, vulnerability severity, and historical exploitability patterns. The primary source of vulnerability metadata is the National Vulnerability Database (NVD), from which we collected CVE records, publication dates, CWE mappings, CVSS v2/v3/v4 scores, and CVE change-history information. To study the evolution of CWE assignments after disclosure, we additionally used the NVD CVE Change History API, which records modifications made to CVE metadata, including the addition and revision of CWE mappings.

To identify vulnerabilities exploited in the wild, we integrated three complementary exploitation datasets: the CISA Known Exploited Vulnerabilities (KEV) catalog\footnote{https://www.cisa.gov/known-exploited-vulnerabilities-catalog}, VulnCheck's exploitation dataset\footnote{https://research.vulncheck.com/}, and  a proprietary dataset of actively exploited CVEs created by the firm GreyNoise\footnote{ https://www.greynoise.io}. A CVE was considered exploited if it appears in at least one of these sources. After merging the datasets and removing duplicate CVE--CWE pairs, the final dataset contained 192,788 unique CVEs, 768 distinct CWE categories, and 3,097 CVEs known to have been exploited in the wild. These exploited CVEs were used to construct historical exploitability profiles for individual CWE categories.

Different research questions require different temporal subsets of the data. RQ1 and RQ2 analyze CVEs published in 2025 to minimize the impact of the well-documented NVD processing backlog during 2024, which delayed vulnerability analysis and metadata publication. RQ3 examines the complete set of available CVE--CWE mappings to characterize the distribution of weakness assignments. RQ4 analyzes temporal exploitability trends using CVEs published between 2020 and 2026, while RQ5 uses the complete merged dataset to identify weakness categories that are disproportionately associated with real-world exploitation. RQ7 and RQ8 analyze CVEs containing multiple CVSS versions to study changes in vulnerability severity across CVSS revisions.

Prior to analysis, duplicate CVE--CWE pairs were removed, and generic or non-informative CWE labels (\texttt{NVD-CWE-noinfo} and \texttt{NVD-CWE-Other}) were excluded where appropriate. Because a single CVE may be associated with multiple CWE identifiers, analyses involving weakness-level exploitability were performed at the CVE--CWE pair level, whereas analyses involving publication dates, CWE assignment delays, and CVSS scores were performed at the CVE level.


\textbf{CWE Association with Exploitation-} To identify software weakness categories that are disproportionately associated with real-world exploitation, as needed in RQ5 and RQ6, we combined CVE--CWE mappings from the National Vulnerability Database (NVD) with exploitation evidence from CISA Known Exploited Vulnerabilities (KEV), VulnCheck, and GrayNoise. A CVE was considered exploited if it appeared in at least one of these exploitation sources. Duplicate CVE--CWE pairs and generic mappings (e.g., \texttt{NVD-CWE-noinfo} and \texttt{NVD-CWE-Other}) were removed prior to analysis. For each CWE, we computed its exploitation rate, contribution to exploited CVEs, and odds ratio measuring the strength of association with exploitation. Statistical significance was assessed using a one-sided Fisher's exact test, and $p$-values were adjusted using the Benjamini--Hochberg procedure to control the false discovery rate. To ensure robust estimates, only CWEs with at least 100 associated CVEs, at least 10 exploited CVEs, observations spanning at least three years, an odds ratio greater than one, and an adjusted $p$-value below 0.05 were retained. We further evaluated the temporal consistency of exploitation using the proportion of years in which a CWE was exploited and a stability score derived from the coefficient of variation of annual exploitation rates. Eligible CWEs were then ranked using a composite score combining exploitation rate (30\%), odds ratio (25\%), exploited-CVE count (25\%), temporal stability (10\%), and the percentage of years with observed exploitation (10\%). Finally, we compared the resulting ranking with the MITRE 2025 CWE Top 25 using set overlap, Jaccard similarity, and Spearman rank correlation, and performed sensitivity analyses under alternative minimum sample-size thresholds to evaluate the robustness of the findings.


The empirical study addresses six aspects of vulnerability analysis. First, we quantify the delay between CVE publication and the initial assignment of a CWE. Second, we examine how frequently CWE mappings are modified after their initial assignment and measure the delay before these updates occur. Third, we characterize the distribution of the number of CWE mappings assigned to individual CVEs. Fourth, we investigate whether historical exploitability rates of CWE categories remain stable across years by comparing annual exploitability profiles and evaluating their predictive relationship over time. Fifth, we identify the CWE categories most strongly associated with vulnerabilities exploited in the wild by combining exploitation likelihood, statistical association, practical contribution, and temporal consistency into a robust ranking and compare the resulting weakness categories with the MITRE 2025 CWE Top 25. Finally, we investigate how CVSS severity scores evolve across CVSS versions and quantify the differences introduced by successive scoring frameworks.

Throughout the empirical study, descriptive statistics are reported using counts, percentages, means, medians, quartiles, and standard deviations where appropriate. Depending on the research question, statistical analyses include Fisher's exact test to evaluate associations between CWE categories and exploitation, Spearman's rank correlation and linear regression to assess temporal stability, Friedman and Wilcoxon signed-rank tests to compare CVSS versions, and bootstrap confidence intervals where applicable. For analyses involving multiple hypothesis tests, $p$-values were adjusted using the Benjamini--Hochberg procedure to control the false discovery rate.

\subsection{Finding 1: Initial CWE Assignments Have Become Faster, but Substantial Delays Persist for Some CVEs}

To motivate the need for automated early CWE assignment, we first examine the delay between the public disclosure of a CVE and the availability of its corresponding CWE mapping. Indeed,  vulnerability prioritization often begins immediately after disclosure. However, CVEs are often published without an assigned CWE, a delaye that hinders practitioners from leveraging weakness-level information during the critical early stages of vulnerability assessment. Quantifying these delays provides insight into the potential benefit of automatically predicting the underlying CWE using an LLM at disclosure time.


\begin{figure}[!t]
    \centering
    \includegraphics[width=\columnwidth]{Figures/figure 5.jpg}
    \caption{Distribution of delays between CVE publication
    and initial CWE assignment across three time periods.}
    \label{fig:cwe_assignment_delay}
\end{figure}



% \begin{figure}[t]
% \centering

% %---------------- Left: Table ----------------%
% \begin{minipage}[c]{0.43\textwidth}
% \vspace{0pt}
% \hspace*{-1cm}
% \centering
% \small

% \textbf{CVE Statistics by Days Until Initial CWE Added}

% \vspace{2mm}

% \resizebox{0.90\linewidth}{!}{%
% \begin{tabular}{lrrr}
% \toprule
% \textbf{Statistic}
% & \textbf{All CVEs}
% & \textbf{2015--2025}
% & \textbf{2025} \\
% \midrule
% Total CVEs & 260{,}982 & 228{,}655 & 46{,}600 \\
% Average    & 65.778    & 24.045    & 12.808 \\
% Median     & 2.236     & 3.166     & 0.000 \\
% Std Dev    & 478.474   & 88.607    & 43.952 \\
% Q1         & 0.000     & 0.000     & 0.000 \\
% Q3         & 7.874     & 8.121     & 0.038 \\
% \bottomrule
% \\
% \end{tabular}%
% }

% \end{minipage}
% \hfill
% %---------------- Right: Figure ----------------%
% \begin{minipage}[t]{0.54\textwidth}
% \vspace{5pt}
% \centering

% \begin{tikzpicture}
% \begin{axis}[
%     width=0.8\linewidth,
%     height=0.65\linewidth,
%     font=\scriptsize,
%     ylabel={Days until initial CWE assignment},
%     ylabel style={font=\scriptsize},
%     xtick={1,2,3},
%     xticklabels={All CVEs, 2015--2025, 2025},
%     xticklabel style={font=\scriptsize},
%     xmin=0.3,
%     xmax=3.7,
%     ymin=-0.5,
%     ymax=16,
%     ymajorgrids=true,
%     grid style={dashed,gray!30},
%     axis x line=bottom,
%     axis y line=left,
%     tick align=outside,
%     boxplot/draw direction=y,
% ]

% % All CVEs
% \addplot+[
%     boxplot prepared={
%         lower whisker=0.0,
%         lower quartile=0.0,
%         median=2.236,
%         upper quartile=7.874,
%         upper whisker=11.811,
%         box extend=0.3,
%     },
%     fill=teal!30,
%     draw=teal!70!black,
% ] coordinates {};

% % 2015--2025
% \addplot+[
%     boxplot prepared={
%         lower whisker=0.0,
%         lower quartile=0.0,
%         median=3.166,
%         upper quartile=8.121,
%         upper whisker=12.182,
%         box extend=0.3,
%     },
%     fill=blue!20,
%     draw=blue!60!black,
% ] coordinates {};

% % 2025
% \addplot+[
%     boxplot prepared={
%         lower whisker=0.0,
%         lower quartile=0.0,
%         median=0.0,
%         upper quartile=0.038,
%         upper whisker=0.057,
%         box extend=0.3,
%     },
%     fill=orange!20,
%     draw=orange!70!black,
% ] coordinates {};

% \end{axis}
% \end{tikzpicture}

% \end{minipage}

% \caption{Summary and distribution of delays between CVE publication and the
% initial CWE assignment.}
% \label{fig:cwe-delay-summary}
% \end{figure}

Figure~\ref{fig:cwe-delay-summary} Shows the distribution of delays for all CVEs, CVEs published between 2015--2025, and CVEs published in 2025. The distributions are highly right-skewed, with most vulnerabilities receiving a CWE assignment within a few days of publication. Across all CVEs, the median delay was 2.24 days, increasing slightly to 3.17 days for CVEs published between 2015 and 2025. In contrast, the median delay for 2025 CVEs was 0 days, indicating that at least half of newly disclosed vulnerabilities received a CWE assignment on the day of publication. Likewise, the upper quartile decreased from approximately eight days in the historical datasets to essentially zero days in 2025, demonstrating substantial improvements in the timeliness of CWE assignment.

Despite these improvements, a small number of vulnerabilities experience exceptionally long assignment delays, producing mean delays of 65.78, 24.05, and 12.81 days for the three datasets, respectively. These long-tail cases substantially inflate the averages and indicate that weakness-level information may remain unavailable for extended periods for some vulnerabilities. Because CWE is one of the features used by exploitability prediction systems such as EPSS, delays in obtaining an official CWE assignment postpone the availability of weakness-level information. Our LLM-based approach addresses this limitation by inferring the underlying CWE immediately after CVE disclosure, enabling earlier vulnerability prioritization using historical weakness-level exploitability information.




\subsection{Finding 2: CWE Assignments Are Rarely Revised, and Most Revisions Occur Shortly After Initial Assignment}

Among the 44,921 CVEs published in 2025, only 905 CVEs (2.01\%) experienced at least one change to their CWE mapping after the initial assignment. These updates accounted for 1,006 CWE change events, indicating that a small number of CVEs underwent multiple mapping revisions. Overall, the time lag between the initial CWE mapping and subsequent updates was very short, with a mean of 0.91 days and a median of 0 days, suggesting that most modifications occurred on the same day as the initial mapping. The distribution of update delays was highly concentrated in the first week: 987 (98.1\%) of all CWE change events occurred within 0–7 days, while only 7 (0.7\%) occurred between 8–30 days and 12 (1.2\%) occurred between 31–90 days. No CWE mapping changes were observed after 90 days. These findings indicate that post-publication changes to CWE mappings are relatively uncommon and, when they do occur, they are typically made shortly after the initial CWE assignment.

We limited our analysis to CVEs published in 2025 to avoid bias introduced by the well-documented NVD processing disruption in 2024. The NVD developed a substantial backlog of vulnerability evaluations starting in February 2024, which caused delays in the assignment and updating of metadata, including CWE mappings. If included in our analysis, these operational delays would have artificially inflated the observed time lag between the initial and subsequent CWE mappings. Restricting the analysis to 2025 CVEs circuments the impact of this backlog and provides a more representative view of normal CWE mapping update behavior ~\cite{NVD2024Backlog,NIST2026NVD}.




\subsection{Finding 3: Most CVEs Are Associated with
a Single CWE, but Missing and Multiple Mappings Persist}

Table~\ref{tab:cwe_mapping_distribution_2020_2025} summarizes the number of CWE mappings assigned to each CVE published between 2020 and 2025. Among the 180,599 CVEs analyzed, the majority (138,343; 76.60\%) were associated with only one CWE, indicating that most vulnerabilities are attributed to a single underlying weakness. An additional 18,763 CVEs (10.39\%) were mapped to two CWEs, while only 844 CVEs (0.47\%) were associated with three CWEs.

A notable proportion of vulnerabilities (22,575; 12.50\%) had no specific CWE mapping, reflecting cases where the underlying weakness had not been identified or was recorded using non-informative classifications such as \texttt{NVD-CWE-noinfo} or \texttt{NVD-CWE-Other}. CVEs associated with four or more CWEs were extremely rare, accounting for fewer than 0.05\% of all vulnerabilities. The maximum number of CWE mappings observed for a single CVE was eight, occurring in only one case.

Overall, these results indicate that while most vulnerabilities can be characterized by a single software weakness, a non-negligible fraction involve multiple underlying weaknesses, suggesting that complex vulnerabilities may span multiple stages of a software flaw or affect several related weakness categories.


% \begin{table}[h]
% \centering
% \caption{Distribution of CVEs by Number of CWEs}
% \begin{tabular}{lrrr}
% \toprule
% \textbf{Number of Assigned CWEs} & \textbf{Number of CVEs} \\
% \midrule
% 1   & 282{,}798    \\
% 2   & 36{,}404       \\
% 3   & 1630        \\
% 4   & 134      \\
% 5   & 20        \\
% 6   & 7        \\
% 7   & 6       \\
% 8   & 1        \\
% \bottomrule
% \end{tabular}
% \end{table}


\begin{table}[t]
\centering
\caption{Distribution of the Number of CWE Mappings per CVE from 2020--2025}
\label{tab:cwe_mapping_distribution_2020_2025}
\begin{tabular}{lrr}
\toprule
\textbf{Number of Mapped CWEs } & \textbf{  Number of CVEs} & \textbf{  Percentage(\%)} \\
\midrule
0 & 22,575  & 12.50 \\
1 & 138,343 & 76.60 \\
2 & 18,763  & 10.39 \\
3 & 844      & 0.47 \\
4 & 57       & 0.03 \\
5 & 10       & 0.01 \\
6 & 3        & 0.00 \\
7 & 3        & 0.00 \\
8 & 1        & 0.00 \\
\midrule
\textbf{Total} & \textbf{180,599} & \textbf{100.00} \\
\bottomrule
\end{tabular}
\end{table}


% \begin{table}[H]
% \centering
% \caption{Distribution of the Number of CWE Mappings per CVEs in 2025}
% \label{tab:cwe_distribution}
% \begin{tabular}{lrr}
% \toprule
% \textbf{Number of CWE Mappings} & \textbf{Number of CVEs} & \textbf{Percentage (\%)} \\
% \midrule
% 0 & 3,528  & 7.85  \\
% 1 & 34,073 & 75.85 \\
% 2 & 6,904  & 15.37 \\
% 3 & 392    & 0.87  \\
% 4 & 18     & 0.04  \\
% 5 & 2      & 0.00  \\
% 6 & 2      & 0.00  \\
% 7 & 1      & 0.00  \\
% 8 & 1      & 0.00  \\
% \bottomrule
% \end{tabular}
% \end{table}




\subsection{Finding 4: Historical CWE Exploitability Exhibits Moderate Temporal Stability}

To evaluate the temporal stability of CWE exploitability profiles, we analyzed CVEs published between 2020 and 2026 and linked them with known exploitation-in-the-wild records from CISA KEV, VulnCheck, and Greynoise's exploitation dataset. To reduce instability caused by sparse observations, we restricted the analysis to CWE-year combinations containing at least 20 CVEs. This resulted in 908 CWE-year observations spanning 210 unique CWEs.

We used spearman rank correlation to compare CWE exploitability rates between consecutive years. All six year-to-year comparisons exhibited statistically significant positive correlations ($p<0.001$), with correlation coefficients ranging from 0.388 (2022--2023) to 0.568 (2021--2022). Across all consecutive years, the mean Spearman correlation was 0.480, indicating moderate temporal stability in CWE exploitability profiles.

To determine whether historical exploitability predicts future exploitability, we compared exploitability rates in year $t$ with those in year $t+1$. The overall lagged Spearman correlation was 0.468 ($p<0.001$), with a bootstrap 95\% confidence interval of [0.396, 0.538], indicating a consistent positive relationship between historical and future exploitability. Linear regression further confirmed this relationship, producing a statistically significant positive slope of 0.364 ($p<0.001$) and an $R^2$ value of 0.184. Although historical exploitability explains only part of the variation in future exploitability, the results demonstrate that CWEs exhibiting higher exploitation rates in one year tend to remain relatively more exploitable in subsequent years.

Finally, we repeated the analysis using only the 76 CWEs that appeared in every year of the study period. The resulting Spearman correlations ranged from 0.450 to 0.641, consistently exceeding those obtained from the complete dataset. This finding suggests that exploitability profiles are even more stable for well-established CWEs that are consistently represented across years.

Overall, these results indicate that CWE exploitability profiles exhibit moderate but statistically significant temporal stability. Historical exploitation patterns therefore provide useful predictive information for estimating the future exploitability of vulnerability classes, although additional factors beyond CWE also influence whether individual vulnerabilities are exploited in the wild.




% \begin{table}[t]
% \centering
% \caption{Summary of Temporal Stability of CWE Exploitability Profiles}
% \label{tab:rq4_summary}
% \begin{tabular}{lr}
% \toprule
% \textbf{Metric} & \textbf{Value} \\
% \midrule
% Study period & 2020--2026 \\
% Unique CWEs & 210 \\
% CWE-year observations & 908 \\
% Minimum CVEs per CWE-year & 20 \\
% Mean consecutive Spearman $\rho$ & 0.480 \\
% Median consecutive Spearman $\rho$ & 0.482 \\
% Lagged Spearman $\rho$ & 0.468 \\
% 95\% CI (Lagged $\rho$) & [0.396, 0.538] \\
% Regression slope & 0.364 \\
% Regression $R^2$ & 0.184 \\
% Regression $p$-value & $<0.001$ \\
% \bottomrule
% \end{tabular}
% \end{table}




\subsection{Finding 5: A Small Set of CWEs Exhibits Persistent Exploitation Associations That Differ from the MITRE Top 25}

\textbf{CWEs Associated with exploitation in the wild-} Among the 192,788 CVEs analyzed, 3,097 were identified as exploited in the wild. To reduce the influence of sparse and potentially unstable estimates, we retained only CWEs associated with at least 100 CVEs and at least 10 exploited CVEs. We further required a positive association with exploitation ($OR>1$) and a Benjamini--Hochberg-adjusted $p$-value below 0.05. Nineteen CWEs satisfied these robustness criteria.

CWE-78 (\emph{Improper Neutralization of Special Elements used in an OS Command}) ranked first according to the composite exploitation score. Of the 4,765 CVEs mapped to CWE-78, 310 were known to be exploited, corresponding to an exploitation rate of 6.51\% and an odds ratio of 4.62. CWE-78 alone accounted for 10.01\% of all exploited CVEs in the dataset and exhibited exploitation in every year of the seven-year study period. CWE-288 (\emph{Authentication Bypass Using an Alternate Path or Channel}) had the highest exploitation rate among the robust categories at 7.87\% and an odds ratio of 5.29, although it accounted for only 1.45\% of all exploited CVEs. Other highly ranked CWEs included CWE-306 (\emph{Missing Authentication for Critical Function}), CWE-434 (\emph{Unrestricted Upload of File with Dangerous Type}), CWE-502 (\emph{Deserialization of Untrusted Data}), CWE-22 (\emph{Path Traversal}), and CWE-287 (\emph{Improper Authentication}).

Several categories combined statistically elevated exploitation likelihood with substantial real-world contribution. CWE-22 (\emph{Path Traversal}) accounted for 6.97\% of exploited CVEs, followed by CWE-862 (\emph{Missing Authorization}) at 5.30\%, CWE-434 (\emph{Unrestricted Upload of File with Dangerous Type}) at 4.84\%, CWE-94 (\emph{Improper Control of Generation of Code ('Code Injection')}) at 4.52\%, CWE-502 (\emph{Deserialization of Untrusted Data}) at 3.78\%, CWE-77 (\emph{Improper Neutralization of Special Elements used in a Command}) at 3.68\%, and CWE-306 (\emph{Missing Authentication for Critical Function}) at 3.52\%. These CWEs were linked to exploitation in each of the seven years of the study, indicating that their observed correlations is not limited to isolated annual events.

Some categories, including CWE-917 (\emph{Improper Neutralization of Special Elements used in an Expression Language Statement ('Expression Language Injection')}), CWE-95 (\emph{Improper Neutralization of Directives in Dynamically Evaluated Code ('Eval Injection')}), and CWE-425 (\emph{Direct Request ('Forced Browsing')}), exhibited relatively high exploitation rates and odds ratios but substantially lower temporal consistency. These CWEs were associated with exploited vulnerabilities in only three or four of the seven study years and had comparatively high coefficients of variation. Therefore, their high aggregate exploitation rates should be interpreted more cautiously than those of CWEs exhibiting persistent exploitation throughout the study period.

\textbf{Comparison with the MITRE CWE Top 25-} Nine of the 19 robust exploitation-associated CWEs also appeared in the MITRE 2025 CWE Top 25: CWE-22 (\emph{Path Traversal}), CWE-77 (\emph{Improper Neutralization of Special Elements used in a Command}), CWE-78 (\emph{Improper Neutralization of Special Elements used in an OS Command}), CWE-94 (\emph{Improper Control of Generation of Code ('Code Injection')}), CWE-306 (\emph{Missing Authentication for Critical Function}), CWE-434 (\emph{Unrestricted Upload of File with Dangerous Type}), CWE-502 (\emph{Deserialization of Untrusted Data}), CWE-862 (\emph{Missing Authorization}), and CWE-918 (\emph{Server-Side Request Forgery (SSRF)}). These shared categories represented 36\% of the MITRE Top 25 and 47.37\% of the robust exploitation-associated set. The Jaccard similarity between the two sets was 0.257. However, the ranks of the nine shared CWEs were not significantly correlated (Spearman's $\rho=-0.067$, $p=0.865$).

The limited ranking agreement reflects the different objectives of the two approaches. The MITRE Top 25 combines weakness prevalence and severity, whereas our ranking focuses specifically on observed exploitation in the wild and incorporates exploitation rate, odds ratio, practical contribution, and temporal consistency. Consequently, some highly prevalent or severe CWEs in the MITRE ranking were not disproportionately associated with observed exploitation, while several CWEs outside the MITRE Top 25 exhibited elevated and persistent exploitation likelihood.

\textbf{Sensitivity Analysis}

We evaluated whether the identified CWE set was sensitive to the minimum sample-size thresholds. Under a relaxed threshold of at least 50 CVEs and five exploited CVEs, all 19 CWEs from the primary analysis remained in the resulting ranking. Increasing the threshold to at least 200 CVEs and 10 exploited CVEs retained 16 of the 19 CWEs (84.21\%), while the strictest threshold of at least 500 CVEs and 20 exploited CVEs retained 15 CWEs (78.95\%). These results indicate that the majority of the identified weakness categories remain associated with exploitation under substantially stricter inclusion requirements, supporting the robustness of the findings.

This finding suggest that vulnerability prioritization should consider historical exploitation characteristics in addition to weakness prevalence and severity. In particular, weaknesses such as CWE-78 (\emph{Improper Neutralization of Special Elements used in an OS Command}), CWE-306 (\emph{Missing Authentication for Critical Function}), CWE-434 (\emph{Unrestricted Upload of File with Dangerous Type}), and CWE-502 (\emph{Deserialization of Untrusted Data}) consistently exhibit elevated exploitation risk across multiple years. Their persistent association with real-world exploitation indicates that these weakness categories should receive high priority during secure software development, vulnerability remediation, and defensive monitoring.




% \begin{table}[H]
% \centering
% \caption{Top 25 CWEs by Mapped Known Exploited Vulnerabilities}
% \begin{tabular}{lrrr}
% \toprule
% \textbf{CWE ID} & \textbf{Mapped KEVs} \\
% \midrule
% 20   & 118    \\
% 78   & 104       \\
% 787   & 100        \\
% 416   & 92      \\
% 119   & 84        \\
% 22   & 75        \\
% 502   & 66       \\
% 94   & 65        \\
% 287   & 39    \\
% 843   & 37       \\
% 306   & 36        \\
% 284   & 33      \\
% 79   & 33        \\
% 264   & 31        \\
% 122   & 30       \\
% 77   & 29        \\
% 89   & 28    \\
% 200   & 24       \\
% 434   & 20        \\
% 918   & 18      \\
% 399   & 18        \\
% 125   & 17        \\
% 288   & 17       \\
% 190   & 16        \\
% 59   & 16        \\
% \bottomrule
% \end{tabular}
% \end{table}

% \begin{table}[H]
% \centering
% \caption{Top 25 CWEs by Mapped Known Exploited Vulnerabilities 2025}
% \begin{tabular}{lrrr}
% \toprule
% \textbf{CWE ID} & \textbf{Mapped KEVs} \\
% \midrule
% 78   & 18    \\
% 502   & 15       \\
% 22   & 13        \\
% 416   & 11      \\
% 787   & 11        \\
% 77   & 7        \\
% 79   & 7       \\
% 287   & 6        \\
% 288   & 6    \\
% 306   & 6       \\
% 94   & 6        \\
% 122   & 5      \\
% 125   & 5        \\
% 121   & 4        \\
% 20   & 4       \\
% 284   & 4        \\
% 36   & 4    \\
% 434   & 4       \\
% 843   & 4        \\
% 862   & 4      \\
% 863   & 4        \\
% 89   & 4        \\
% 119   & 3       \\
% 120   & 3        \\
% 506   & 3        \\
% \bottomrule
% \end{tabular}
% \end{table}



% % \textbf{RQ6}: Can historical CWE exploitability improve the prioritization of newly disclosed CVEs? 
% % This research question evaluates whether incorporating historical exploitability information at the CWE level provides better prioritization of newly disclosed CVEs than using severity-based metrics alone.

% % Methodology: Compute the historical exploitability score for a CWE using only CVEs disclosed before the target CVE. Evaluate CVEs disclosed in 2025 and 2024, its CWE score should be based only on cves before their release date. Extract their CWE exploitability score and check if they really exploitated in the wild after their release. Check if there is any meaningful correlation.

% % if cve-x matched to CWE-y, and for all cves released before the date of cve-x and matched to CWE-y, get their exploitation in the wild. we have exploitation for cve-x in the wild as well. 

% % high correlation or any statistic? 
% % statistical test?

% RQ6: Does incorporating historical CWE exploitability improve the ranking of newly disclosed CVEs for exploitation prioritization compared with CVSS severity alone?








% \subsection{\textbf{RQ6}: How often do vulnerabilities receive lower CVSS scores under newer CVSS versions, and what is the magnitude of these score reductions?}

% As shown in Table~\ref{tab:cvss_score_reduction}, vulnerabilities are substantially more likely to receive lower severity scores under CVSS v4 than under CVSS v3. Among the 5,919 CVEs with both CVSS v2 and v3 scores, 1,111 (18.77\%) received a lower score under CVSS v3. These reductions were generally small, with a mean decrease of 0.39 points, a median of 0.2 points, and a maximum reduction of 2.5 points. In contrast, among the 10,509 CVEs with both CVSS v3 and v4 scores, 7,963 (75.77\%) received a lower score under CVSS v4. The magnitude of these reductions was much larger, with a mean decrease of 2.83 points, a median of 2.9 points, and a maximum reduction of 9.8 points. Overall, the findings indicate that score decreases are relatively uncommon and modest when moving from CVSS v2 to v3, but highly frequent and substantial when moving from CVSS v3 to v4.

% \begin{table}[t]
% \centering
% \caption{Comparison of CVSS Score Reductions Across Versions for CVEs Published in 2025}
% \label{tab:cvss_score_reduction}
% \begin{tabular}{lcc}
% \toprule
% \textbf{Comparison} & \textbf{\% Lower in New Version} & \textbf{Mean Reduction} \\
% \midrule
% CVSS v3 vs. CVSS v2 & 18.77\% & 0.39 \\
% CVSS v4 vs. CVSS v3 & 75.77\% & 2.83 \\
% \bottomrule
% \end{tabular}
% \end{table}

% The substantial reduction in CVSS v4 scores is consistent with the design objectives of the CVSS v4.0 specification. Unlike CVSS v3.1, CVSS v4 introduces the Attack Requirements (AT) metric, refines the calculation of vulnerability impacts, and adopts a redesigned scoring formula that more precisely models exploitability and impact. The specification also separates intrinsic vulnerability severity from contextual factors through enhanced Threat, Environmental, and Supplemental metrics, reducing the reliance on the Base Score alone for prioritization. Collectively, these changes reduce score inflation observed in CVSS v3.1 and produce a broader distribution of Base Scores, explaining why a large proportion of vulnerabilities receive lower severity scores under CVSS v4.0.



% % \subsection{\textbf{RQ7}: How significant are the differences in vulnerability severity scores across CVSS versions?}

% % A Friedman test was conducted to evaluate whether vulnerability severity scores differ across CVSS versions. The results revealed a statistically significant difference among CVSS v2, v3, and v4 scores ($\chi^2 = 6379.49$, $p < 0.001$), indicating that the three versions do not assign equivalent severity scores. Pairwise Wilcoxon signed-rank tests with Bonferroni correction further confirmed statistically significant differences between both version pairs ($p < 0.001$), as summarized in Table~\ref{tab:cvss_version_comparison}. Compared with CVSS v2, CVSS v3 assigned higher severity scores to 4,773 (80.64\%) of the 5,919 comparable vulnerabilities, lower scores to 1,111 (18.77\%), and identical scores to 35 (0.59\%). The mean severity score increased from 6.19 to 7.83, corresponding to a large positive effect size (rank-biserial correlation $r = 0.903$; Cohen's $d = 1.246$). In contrast, compared with CVSS v3, CVSS v4 assigned lower severity scores to 7,963 (75.77\%) of the 10,509 comparable vulnerabilities, higher scores to 2,431 (23.13\%), and identical scores to 115 (1.09\%). The mean severity score decreased from 7.51 to 5.61, with a large negative effect size (rank-biserial correlation $r = -0.714$; Cohen's $d = -0.773$). These findings indicate that CVSS v3 systematically increased vulnerability severity relative to CVSS v2, whereas CVSS v4 systematically reduced severity scores relative to CVSS v3.

% % \begin{table}[t]
% % \centering
% % \caption{Comparison of Vulnerability Severity Scores Across CVSS Versions}
% % \label{tab:cvss_version_comparison}
% % \begin{tabular}{lrrrrrr}
% % \toprule
% % \textbf{Comparison} & \textbf{N} & \textbf{\% Higher} & \textbf{\% Lower} & \textbf{$r$} & \textbf{$d$} & \textbf{$p$-value} \\
% % \midrule
% % CVSS v2 vs. CVSS v3 &
% % 5,919 &
% % 80.64 &
% % 18.77 &
% % 0.903 &
% % 1.246 &
% % $<0.001$ \\

% % CVSS v3 vs. CVSS v4 &
% % 10,509 &
% % 23.13 &
% % 75.77 &
% % $-0.714$ &
% % $-0.773$ &
% % $<0.001$ \\
% % \bottomrule
% % \end{tabular}
% % \end{table}





% \subsection{Discussion}


% Our findings in RQ4 provide empirical justification for incorporating historical CWE exploitability into CVE risk assessment. Although exploitability is not completely determined by the weakness type, CWE exploitability profiles exhibit moderate and statistically significant temporal stability across multiple years. Furthermore, historical exploitability rates are positively correlated with future exploitability ($\rho=0.468$, $p<0.001$), indicating that CWEs with consistently higher exploitation rates tend to remain relatively more attractive to attackers over time. These results suggest that the historical exploitability of a CWE constitutes a meaningful predictive feature for estimating the likelihood that newly disclosed CVEs associated with that CWE will be exploited in the wild. Nevertheless, the observed $R^2$ value of 0.184 indicates that CWE history alone is insufficient for accurate prediction and should be combined with additional vulnerability characteristics, such as CVSS metrics, affected products, and exploit availability, to improve predictive performance.



\subsection{Discussion}
\label{sec:empirical_discussion}

The empirical findings provide the motivation and
supporting evidence for incorporating historical
CWE exploitability into early vulnerability
prioritization. Our analysis highlights two key
aspects of this approach: the availability and
reliability of CWE assignments, and the extent
to which historical exploitation patterns can
inform the assessment of newly disclosed
vulnerabilities.

First, although the timeliness of initial CWE
assignments has improved, some vulnerabilities
continue to experience substantial delays.
While the median assignment delay for CVEs
published in 2025 was zero days, the mean
remained 12.81 days, reflecting a long-tailed
distribution. Consequently, weakness-level
information may not be available when a
vulnerability is first disclosed. This finding
motivates automated CWE prediction to enable
the use of weakness-level information before
an official assignment becomes available.

Second, CWE mappings are generally stable after
their initial assignment. Only 2.01\% of CVEs
published in 2025 experienced subsequent CWE
mapping changes, with most revisions occurring
within the first week. Furthermore, 76.60\%
of CVEs published between 2020 and 2025 were
associated with exactly one CWE. These findings
suggest that single-label CWE prediction is a
reasonable starting point for automated mapping,
although missing and multiple CWE assignments
must also be considered. Importantly, the
stability of official assignments does not
guarantee the accuracy of automated predictions,
which must be evaluated independently.

More importantly, our analysis demonstrates
that historical CWE exploitability exhibits
moderate temporal stability. Exploitation rates
were positively correlated across consecutive
years, with an overall lagged Spearman correlation
of $\rho=0.468$ ($p<0.001$). This indicates
that CWE categories with relatively high
exploitation rates tend to maintain higher
rates in subsequent years. In addition, we
identified 19 CWE categories with statistically
significant and robust associations with
exploitation in the wild. Several of these
weaknesses, including command injection,
authentication-related weaknesses, unrestricted
file upload, deserialization, and path traversal,
exhibited persistent exploitation across
multiple years.

The limited agreement between our exploitation-
based CWE ranking and the MITRE 2025 CWE Top 25
further highlights the importance of considering
observed exploitation patterns. While the MITRE
ranking emphasizes weakness prevalence and
severity, our analysis identifies categories
associated with elevated historical exploitation
likelihood. These perspectives are complementary,
but they serve different vulnerability management
objectives.

Taken together, these findings support the use
of historical CWE exploitation rates as an
informative signal for early vulnerability
prioritization. Although historical CWE rates
explain only part of the variation in subsequent
annual exploitation rates ($R^2=0.184$),
their statistically significant temporal
association suggests that they can provide
useful information when vulnerability-specific
exploitation evidence is limited or unavailable.
The objective is therefore not to determine
exploitation solely from the CWE identifier,
but to use historical weakness-level knowledge
to support the initial ranking of newly
disclosed vulnerabilities.

These observations establish the motivation
for the next component of our framework:
automatically identifying the underlying CWE
of a newly disclosed CVE from its description.
By linking the predicted CWE to its historical
exploitability profile, the proposed approach
can assign an initial exploitation-based
priority without waiting for an official CWE
mapping. The subsequent sections investigate
the accuracy of automated CWE identification
and the effectiveness of historical CWE
information for vulnerability prioritization.



\section{CWE-based CVE Exploitability Estimation Experiment}

The objective of our framework is to automatically map a newly disclosed CVE to its underlying Common Weakness Enumeration (CWE) using a Large Language Model (LLM), thereby enabling exploitability estimation before real-world exploitation evidence becomes available. As demonstrated in the previous section, different CWE categories exhibit substantially different historical exploitation characteristics. Therefore, once a CVE is accurately mapped to its corresponding CWE, the historical exploitability profile associated with that weakness can be used to estimate the likelihood that the new vulnerability will be exploited in the wild. Consequently, the accuracy of CWE classification directly determines the reliability of the proposed exploitability estimation framework.





\subsection{Why Not Use Generative LLMs for CWE Mapping?}
\label{subsec:generative_llm_limitations}

Generative large language models (LLMs), such as ChatGPT,
offer a straightforward approach to automated CVE-to-CWE
mapping: given a vulnerability description, the model
can be prompted to identify the corresponding software
weakness. However, CWE classification requires selecting
a precise identifier from a large, hierarchical taxonomy
containing semantically overlapping weaknesses.
Generative models may confuse CWE identifiers and
definitions, select overly general weakness categories,
or produce inconsistent predictions. To investigate
whether generative LLMs provide sufficiently reliable
CWE assignments for our exploitability estimation
framework, we conducted a preliminary experimental
evaluation of several locally deployed models.

We evaluated Llama 3.2 (3B), Qwen 2.5 (14B),
Qwen 2.5-Coder (14B), and Phi-4 (14B) using CVEs
published in January 2024, with NVD CWE assignments
as reference labels. Preliminary experiments showed
that allowing LLMs to generate CWE identifiers
without an explicit candidate list frequently
resulted in confusion between CWE identifiers,
names, and descriptions. We therefore developed
a constrained, hierarchy-aware prompting pipeline
that supplied official CWE identifiers, titles,
and descriptions. Three configurations were
investigated: (A) classification using 41 top-level
CWE-699 categories, (B) expansion to 157
software-related categories, and (C) top-three
parent selection followed by classification
among their combined child weaknesses.

Initial experiments on 25 CVEs demonstrated that
expanding the taxonomy and introducing multiple
parent candidates improved the performance of
Qwen 2.5, increasing exact-match accuracy from
4.0\% to 20.0\%. However, the smaller Llama 3.2
model exhibited prediction collapse under
Configuration B, repeatedly selecting the same
CWE for 24 of 25 vulnerabilities. This behavior
illustrates that apparent classification accuracy
may conceal limited discrimination among weaknesses.

We subsequently evaluated Configuration C on
150 CVEs. At temperature $T=0$, Qwen 2.5-Coder
achieved the highest exact-match accuracy (23.3\%)
and parent-level accuracy (28.0\%), outperforming
Qwen 2.5 (17.3\% and 20.7\%), Phi-4 (16.7\%
and 16.0\%), and Llama 3.2 (7.3\% and 9.3\%).
Adjusting the decoding temperature improved
Qwen 2.5-Coder's exact-match accuracy to 26.0\%
at $T=0.3$. Furthermore, enriching the CVE
description with information retrieved from
external reference links increased exact-match
accuracy to 32.0\% and parent-level accuracy
to 32.7\%, although average processing time
increased from 3.9 to 6.4 seconds per CVE.

Despite these refinements, the best-performing
generative configuration correctly identified
the exact CWE for only 32.0\% of the evaluated
vulnerabilities. Moreover, 10.7\% of the
reference CWE assignments were outside the
evaluated taxonomy, illustrating the importance
of candidate coverage. These findings suggest
that prompt engineering, hierarchical reasoning,
and contextual enrichment alone do not guarantee
reliable fine-grained CWE identification.
The observed sensitivity to model selection,
taxonomy construction, and decoding parameters
also complicates the use of generative models
as a consistent automated classification mechanism.

These limitations motivate our investigation of
a lightweight, fine-tuned bi-encoder as an
alternative to generative CWE prediction.
Rather than generating CWE identifiers through
autoregressive decoding, the proposed approach
learns semantic representations of CVE descriptions
and CWE definitions and ranks predefined candidate
weaknesses according to their similarity.
This design avoids unconstrained identifier
generation and provides a deterministic
candidate-ranking mechanism. The following
section introduces the proposed bi-encoder
architecture and evaluates whether task-specific,
parameter-efficient fine-tuning can improve
CWE identification.


\begin{table}[t]
\centering
\caption{Generative LLM performance for CVE-to-CWE
mapping ($n=150$). Results use Configuration C
and $T=0$, unless otherwise specified.}
\label{tab:generative_llm_summary}
\scriptsize
\setlength{\tabcolsep}{3pt}
\begin{tabular}{lrr}
\toprule
\textbf{Model / Configuration}
& \textbf{Exact (\%)}
& \textbf{Parent (\%)} \\
\midrule
Llama 3.2 (3B)        & 7.3  & 9.3 \\
Qwen 2.5 (14B)        & 17.3 & 20.7 \\
Phi-4 (14B)          & 16.7 & 16.0 \\
Qwen 2.5-Coder (14B)  & 23.3 & 28.0 \\
\midrule
Qwen 2.5-Coder ($T=0.3$)
                       & 26.0 & 26.7 \\
+ Reference enrichment & 32.0 & 32.7 \\
\bottomrule
\end{tabular}
\end{table}


It is important to note that the evaluated generative
LLMs may have encountered some of the CVE descriptions
and their corresponding CWE assignments during
pretraining, as this information is publicly available
through vulnerability databases. Although the extent
of such exposure cannot be verified, the possibility
of training-data overlap makes the relatively low
classification accuracy particularly noteworthy.
Even with potential prior exposure to some CVE--CWE
mappings, the best-performing configuration achieved
only 32.0\% exact-match accuracy. This observation
further motivates the development of a task-specific
CWE classification approach.

To address this concern in our proposed bi-encoder
evaluation, we adopt a temporal data-splitting
strategy, using CVEs published between 1999 and
2021 for training and CVEs published between 2023
and 2025 for testing. This separation prevents
direct overlap between the CVE records used for
fine-tuning and those used for evaluation, providing
a more realistic assessment of the model's ability
to generalize to newly disclosed vulnerabilities.

To support reproducibility, all prompts used in the
generative LLM experiments, along with the implementation
and evaluation scripts, are publicly available in our
repository at
\url{https://anonymous.4open.science/r/cwe-classifier-337D/}.


% \subsection{LLM based CWE mapping Overview}

% The proposed classification pipeline leverages the hierarchical organization of the CWE taxonomy instead of requiring the LLM to select directly from hundreds of weakness entries. According to the MITRE CWE Usage Guidance~\cite{mitre_cwe_guidance}, weaknesses are organized into multiple abstraction levels, including \emph{Pillar}, \emph{Class}, \emph{Base}, and \emph{Variant}. Pillar entries represent broad conceptual categories, Class entries group related weaknesses, Base entries describe specific weakness types appropriate for vulnerability classification, and Variant entries provide technology- or implementation-specific refinements. This hierarchical organization naturally supports progressive refinement, allowing the search space to be reduced while preserving semantic relationships among weaknesses.

% Our pipeline adopts this hierarchy as a coarse-to-fine classification strategy. Given the textual description of a CVE, the LLM first identifies the most relevant high-level weakness category from a predefined set of Pillar and Class candidates. The candidate space is then restricted to the corresponding child weaknesses, from which the model selects the most appropriate specific CWE. Each candidate is represented using its official CWE identifier, title, and MITRE description to improve semantic matching and reduce ambiguity.

% To investigate how taxonomy design and search strategy affect classification performance, we evaluate three progressively refined pipeline configurations that differ in taxonomy coverage and hierarchical search strategy. These configurations are described in the following subsection.

% To support reproducibility, we make the implementation of the proposed LLM-based CWE classification framework, including the prompting pipeline, taxonomy construction, and evaluation scripts, publicly available through an anonymous repository at \url{https://anonymous.4open.science/r/cwe-classifier-337D/}.

% \subsection{Experimental Setup}
% % number of cves?1000, all those that are have CWE mapped in https://cwe.mitre.org/data/definitions/699.html
% %cve->NVD mapped CWE
% % how choose them? randomly 2025-2021
% % how you check they are related to software? nvd related to software

% % step 1:maped to Pillar
% % step 2:maped to base/class/varient/



% The objective of the experimental evaluation is to determine how accurately an LLM can infer the underlying software weakness from a CVE description, since the reliability of the subsequent exploitability estimation depends directly on the correctness of this CWE assignment. To investigate the effects of taxonomy coverage and hierarchical search strategy, we evaluated three progressively refined pipeline configurations. Each configuration builds upon the previous one by expanding the taxonomy or reducing hierarchical error propagation.


% \begin{itemize}
%   \item \textbf{Configuration A (single-parent, baseline taxonomy).} A two-step
%   pipeline. Step 1 gives the model the CVE description plus the 41
%   top-level CWE-699 (``Software Development'') categories and asks for the
%   single best-fitting parent. Step 2 gives the model the CVE description
%   plus that parent's enumerated children and asks for the single
%   best-fitting specific CWE-ID.

%   \item \textbf{Config B (single-parent, expanded software taxonomy).}
%   Identical to Configuration A, but the Step 1 candidate list is augmented with
%   116 additional software-relevant Pillar- and Class-level weaknesses drawn
%   from the CWE-1000 (``Research Concepts'') view (e.g., CWE-119, CWE-200,
%   CWE-284, CWE-862, CWE-269, CWE-311) -- broader IDs that NVD analysts
%   frequently assign directly but which are not represented as leaves
%   anywhere in the CWE-699 view. Weaknesses exclusive to the CWE-1194
%   view (``Hardware Design'') were excluded to keep the taxonomy scoped to
%   software, raising the candidate pool from 41 to 157 top-level
%   categories.

%   \item \textbf{Configuration C (top-$K$ parent branching, - expanded taxonomy).}
%   % CWE list similar to B
% % evaluation, if the NVD mapped CWE is in 3 the class or pillar of the LLM mapped CWE, LLM identify it correctly!

%   Addresses error-compounding in Configurations A/B, where an incorrect Step 1
%   parent choice made the correct Step 2 leaf permanently unreachable. Step
%   1 requests the top 3 ranked candidate parents instead of one; Step 2
%   pools the enumerated children of all 3 chosen parents (deduplicated),
%   including parents with no enumerated children as leaf candidates in
%   their own right, and always performs an explicit arbitration pass over
%   the pooled set.
  
% \end{itemize}

% % All candidate lists included each CWE's official name and description in
% % every configuration, so that differences in results reflect changes to
% % taxonomy coverage and pipeline architecture rather than information
% % availability. Importantly, early tests of allowing local LLM models to 
% % freely select CWE's from their own training was found to be extremely inaccurate 
% % as CWE numbers, names, and descriptions would often be confused, meaning sample 
% % lists would need to be provided to the models.
% % Models were served locally via Ollama on an NVIDIA RTX 3080
% % Ti (16GB VRAM). Two models were evaluated at the $n=25$ scale --
% % llama3.2 (3B, Q4\_K\_M) and qwen2.5 (14B, Q4\_K\_M) -- both at temperature
% % $=0$ (greedy decoding). Ground truth was 25 CVE records randomly sampled
% % from CVEs published between 2024-01-01 and 2024-01-03, retrieved from the
% % NVD CVE API and restricted to records carrying at least one substantive
% % CWE assignment (placeholder labels such as \texttt{NVD-CWE-noinfo} and
% % \texttt{NVD-CWE-Other} were excluded). The same 25 CVEs were used across
% % every model/configuration pair to allow direct, paired comparison.
% All candidate lists included the official CWE identifier, title, and description so that differences in performance reflected only changes in taxonomy coverage and pipeline architecture rather than information availability. Preliminary experiments allowing local LLMs to freely select CWE identifiers from their internal knowledge produced consistently poor results, as models frequently confused CWE identifiers, titles, and descriptions. Consequently, every experiment explicitly supplied the candidate CWE list as part of the prompt.

% Models were executed locally through Ollama on an NVIDIA RTX 3080 Ti (16 GB VRAM). Initial experiments evaluated llama3.2 (3B, Q4\_K\_M) and qwen2.5 (14B, Q4\_K\_M) using deterministic decoding (temperature $=0$). Ground truth consisted of 25 randomly sampled CVEs published between January 1 and January 3, 2024, obtained from the NVD API and restricted to vulnerabilities with substantive CWE assignments after excluding placeholder values \texttt{NVD-CWE-noinfo} and \texttt{NVD-CWE-Other}. The identical evaluation set was used across all configurations to enable paired comparisons.

% \subsection{Evaluation Metrics}
% The quality of the predicted CWE assignments was evaluated using the following metrics:

% \begin{itemize}
%     \item \textbf{Exact-match accuracy:} Measures whether the predicted CWE identifier exactly matches at least one of the CWE identifiers assigned to the CVE by the NVD.

%     \item \textbf{Parent-level accuracy:} Measures whether the predicted CWE and at least one ground-truth CWE belong to the same high-level parent category in the CWE hierarchy, providing partial credit for semantically related predictions.

%     \item \textbf{Taxonomy coverage:} Measures the proportion of ground-truth CWE identifiers that are represented within the evaluated taxonomy. This metric ensures that classification errors are not attributed to missing candidate weaknesses.

%     \item \textbf{Prediction diversity:} Measures the number of distinct CWE identifiers predicted across the evaluation dataset. This metric complements accuracy by identifying degenerate prediction behavior, where a model repeatedly predicts the same common weakness regardless of the input.
% \end{itemize}

% The combination of these metrics provides a comprehensive assessment of the proposed CWE classification pipeline by evaluating classification correctness, hierarchical consistency, taxonomy completeness, and prediction robustness. Since the estimated exploitability of a CVE depends directly on its predicted CWE, improvements in these metrics translate into more reliable exploitability estimation.


% % \begin{itemize}
% %   \item \textbf{Exact-match accuracy}: the predicted CWE-ID matches at
% %   least one of the CWE-IDs officially assigned to the CVE by NVD.
% %   \item \textbf{Parent-level accuracy}: the predicted CWE-ID falls under
% %   the same top-level category as at least one ground-truth CWE-ID.
% %   \item \textbf{Taxonomy coverage}: the proportion of ground-truth CWE-IDs
% %   that exist within the evaluated taxonomy.
% %   \item \textbf{Prediction diversity}: the number of distinct CWE-IDs
% %   predicted across a sample, out of the sample size. This metric was
% %   introduced after Config B revealed that raw accuracy alone can be
% %   misleading -- a model that repeats a single prediction regardless of
% %   input can still score well if that prediction coincides with several
% %   ground-truth labels in the sample.
% % \end{itemize}

% \subsection{Pipeline Configuration Comparison ($n=25$)}
% We first evaluate how taxonomy coverage and hierarchical search strategy influence CWE classification accuracy. Table~\ref{tab:config_n25} summarizes exact-match accuracy, parent-level accuracy, and prediction diversity for the two evaluated models under each pipeline configuration.


% Config B's headline numbers for llama3.2 are misleading:
% the model repeated a single prediction 24 of 25 times, so its 12\%/20\%
% scores reflect that one repeated guess coinciding with several
% ground-truth labels rather than genuine discrimination. Config C corrects
% for this -- llama3.2's scores drop to 8\%/8\% but predictions become
% meaningfully more varied (11/25 unique), a more honest reflection of its
% limited capability. For qwen2.5 (14B), each configuration change produced
% a genuine, diversity-preserving improvement, culminating in Config C's
% 20.0\% exact-match and 20.0\% parent-level accuracy with 17 of 25
% predictions distinct.

% \vspace{-.7cm}
% \begin{table}[ht]
% \centering
% \caption{Exact-match and parent-level accuracy by pipeline configuration (n = 25).}
% \label{tab:config_n25}
% \begin{tabular}{lrrcrrc}
% \toprule
% & \multicolumn{3}{c}{\textbf{llama3.2 (3B)}} & \multicolumn{3}{c}{\textbf{qwen2.5 (14B)}} \\
% \cmidrule(lr){2-4} \cmidrule(lr){5-7}
% \textbf{Configuration} & \textbf{Exact} & \textbf{ Parent} & \textbf{ Uniq} & \textbf{Exact} & \textbf{ Parent} & \textbf{ Uniq} \\
% \midrule
% Config A (41 cat., single parent)    & 0.0\%  & 0.0\%  & n/a   & 4.0\%  & 16.0\% & n/a   \\
% Config B (157 cat., single parent)   & 12.0\% & 20.0\% & 1/25  & 16.0\% & 16.0\% & 20/25 \\
% Config C (157 cat., top-3 branching) & 8.0\%  & 8.0\%  & 11/25 & 20.0\% & 20.0\% & 17/25 \\
% \bottomrule
% \end{tabular}
% \end{table}

% \vspace{-.5cm}
% \subsection{Large-Sample Model Comparison ($n=150$)}
% After identifying the most effective pipeline architecture, we investigated whether the observed trends generalize to a larger evaluation dataset and whether domain-specialized language models outperform general-purpose models. Configuration C was therefore evaluated on an independent sample of 150 CVEs drawn from 1,195 eligible vulnerabilities published between January 1 and January 15, 2024. In addition to llama3.2 (3B) and qwen2.5 (14B), two additional 14B models were evaluated: qwen2.5-coder, a code-specialized variant, and phi4, a reasoning-oriented model.


% Qwen2.5-coder (14B) produced the best result of the entire study by a
% clear margin -- roughly 6 percentage points above its general-purpose
% sibling on both exact-match and parent-level accuracy, at essentially the
% same inference cost (3.9s vs.\ 3.8s per CVE). phi4 (14B) did not
% outperform the general-purpose qwen2.5 (14B) baseline despite its
% reasoning-oriented training: comparable exact-match, slightly worse
% parent-level accuracy, and the slowest inference of the four models
% (5.6s per CVE), partly because its raw replies occasionally included
% explanatory prose despite an explicit ``no other text'' instruction,
% requiring a fallback extraction step. All three 14B-class models
% comfortably outperformed llama3.2 (3B), consistent with the earlier
% scale comparison.

% \begin{table}[ht]
% \centering
% \caption{Model comparison under Config C (n=150). Ground truth outside the taxonomy was 16/150 (10.7\%), identical across all four models.}
% \label{tab:model_n150}
% \begin{tabular}{lrrrr}
% \toprule
% \textbf{Metric} & \textbf{llama3.2 } & \textbf{ qwen2.5 } & \textbf{ qwen2.5-coder} & \textbf{phi4} \\
% \midrule
% Exact-match accuracy  & 7.3\%  & 17.3\% & 23.3\% & 16.7\% \\
% Parent-level accuracy & 9.3\%  & 20.7\% & 28.0\% & 16.0\% \\
% Prediction diversity  & 30/150 & 47/150 & 38/150 & 34/150 \\
% Avg.\ time per CVE    & 1.2s   & 3.8s   & 3.9s   & 5.6s   \\
% \bottomrule
% \end{tabular}
% \end{table}

% \subsection{Temperature Sensitivity}The previous experiments used deterministic decoding (temperature $=0$), which is commonly employed for structured prediction tasks. We next investigated whether introducing controlled randomness improves CWE classification performance. The best-performing model identified in the previous experiment, qwen2.5-coder (14B), was evaluated using the same 150-CVE dataset under four temperature settings.

% Accuracy was not maximized at temperature $=0$. Both moderate settings
% (0.3 and 0.7) outperformed greedy decoding on exact-match, and temperature
% 0.7 produced the best parent-level accuracy of any run in the study
% (30.7\%). Prediction diversity rose steadily with temperature, but this
% relationship decoupled from accuracy at temperature $=1.0$: diversity
% peaked (57/150 unique) while both accuracy metrics fell below the
% temperature $=0$ baseline. This is consistent with mild randomness helping
% the model escape a locally-favored but incorrect greedy token choice,
% while excessive randomness begins injecting unwarranted guesses.

% \begin{table*}[ht]
% \centering
% \caption{Temperature sensitivity of the evaluated LLMs under Config C ($n=150$).}
% \label{tab:temperature_all}
% \small
% \begin{tabular}{llcccc}
% \toprule
% \textbf{Model} & \textbf{Metric} &
% \textbf{T=0.0} &
% \textbf{T=0.3} &
% \textbf{T=0.7} &
% \textbf{T=1.0} \\
% \midrule

% \multirow{4}{*}{qwen2.5-coder (14B)}
% & Exact-match accuracy  & 23.3\% & 26.0\% & 25.3\% & 19.3\% \\
% & Parent-level accuracy & 28.0\% & 26.7\% & 30.7\% & 22.7\% \\
% & Prediction diversity  & 38/150 & 45/150 & 46/150 & 57/150 \\
% & Avg.\ time (s/CVE)    & 3.9 & 3.9 & 3.8 & 3.8 \\
% \midrule

% \multirow{4}{*}{qwen2.5 (14B)}
% & Exact-match accuracy  & 17.3\% & 14.0\% & 18.0\% & 16.7\% \\
% & Parent-level accuracy & 20.7\% & 16.7\% & 21.3\% & 20.0\% \\
% & Prediction diversity  & 47/150 & 47/150 & 48/150 & 55/150 \\
% & Avg.\ time (s/CVE)    & 3.8 & 3.7 & 3.6 & 3.6 \\
% \midrule

% \multirow{4}{*}{phi4 (14B)}
% & Exact-match accuracy  & 16.7\% & 16.7\% & 16.7\% & 16.7\% \\
% & Parent-level accuracy & 16.0\% & 18.0\% & 18.0\% & 16.7\% \\
% & Prediction diversity  & 34/150 & 42/150 & 51/150 & 52/150 \\
% & Avg.\ time (s/CVE)    & 5.6 & 5.3 & 5.3 & 5.4 \\
% \midrule

% \multirow{4}{*}{llama3.2 (3B)}
% & Exact-match accuracy  & 7.3\% & 5.3\% & 3.3\% & 2.0\% \\
% & Parent-level accuracy & 9.3\% & 6.7\% & 4.7\% & 2.0\% \\
% & Prediction diversity  & 30/150 & 38/150 & 64/150 & 78/150 \\
% & Avg.\ time (s/CVE)    & 1.2 & 1.2 & 1.2 & 1.1 \\
% \bottomrule
% \end{tabular}
% \end{table*}


% % \begin{table}[ht]
% % \centering
% % \caption{Temperature sensitivity, qwen2.5-coder (14B), Config C (n=150).}
% % \label{tab:temperature-qwen25coder}
% % \begin{tabular}{lrrrr}
% % \toprule
% % \textbf{Metric} & \textbf{ T=0.0 } & \textbf{ T=0.3 } & \textbf{ T=0.7 } & \textbf{ T=1.0 } \\
% % \midrule
% % Exact-match accuracy  & 23.3\% & 26.0\% & 25.3\% & 19.3\% \\
% % Parent-level accuracy & 28.0\% & 26.7\% & 30.7\% & 22.7\% \\
% % Prediction diversity  & 38/150 & 45/150 & 46/150 & 57/150 \\
% % Avg.\ time per CVE    & 3.9s   & 3.9s   & 3.8s   & 3.8s   \\
% % \bottomrule
% % \end{tabular}
% % \end{table}

% % \begin{table}[ht]
% % \centering
% % \caption{Temperature sensitivity, qwen2.5 (14B), Config C (n=150).}
% % \label{tab:temperature-qwen25}
% % \begin{tabular}{lrrrr}
% % \toprule
% % \textbf{Metric} & \textbf{ T=0.0 } & \textbf{ T=0.3 } & \textbf{ T=0.7 } & \textbf{ T=1.0 } \\
% % \midrule
% % Exact-match accuracy  & 17.3\% & 14.0\% & 18.0\% & 16.7\% \\
% % Parent-level accuracy & 20.7\% & 16.7\% & 21.3\% & 20.0\% \\
% % Prediction diversity  & 47/150 & 47/150 & 48/150 & 55/150 \\
% % Avg.\ time per CVE    & 3.8s   & 3.7s   & 3.6s   & 3.6s   \\
% % \bottomrule
% % \end{tabular}
% % \end{table}


% % \begin{table}[ht]
% % \centering
% % \caption{Temperature sensitivity, phi4 (14B), Config C (n=150).}
% % \label{tab:temperature-phi4}
% % \begin{tabular}{lrrrr}
% % \toprule
% % \textbf{Metric} & \textbf{ T=0.0 } & \textbf{ T=0.3 } & \textbf{ T=0.7 } & \textbf{ T=1.0 } \\
% % \midrule
% % Exact-match accuracy  & 16.7\% & 16.7\% & 16.7\% & 16.7\% \\
% % Parent-level accuracy & 16.0\% & 18.0\% & 18.0\% & 16.7\% \\
% % Prediction diversity  & 34/150 & 42/150 & 51/150 & 52/150 \\
% % Avg.\ time per CVE    & 5.6s   & 5.3s   & 5.3s   & 5.4s   \\
% % \bottomrule
% % \end{tabular}
% % \end{table}

% % \begin{table}[ht]
% % \centering
% % \caption{Temperature sensitivity, llama3.2 (3B), Config C (n=150).}
% % \label{tab:temperature-llama32}
% % \begin{tabular}{lrrrr}
% % \toprule
% % \textbf{Metric} & \textbf{ T=0.0 } & \textbf{ T=0.3 } & \textbf{ T=0.7 } & \textbf{ T=1.0 } \\
% % \midrule
% % Exact-match accuracy  & 7.3\%  & 5.3\%  & 3.3\%  & 2.0\%  \\
% % Parent-level accuracy & 9.3\%  & 6.7\%  & 4.7\%  & 2.0\%  \\
% % Prediction diversity  & 30/150 & 38/150 & 64/150 & 78/150 \\
% % Avg.\ time per CVE    & 1.2s   & 1.2s   & 1.2s   & 1.1s   \\
% % \bottomrule
% % \end{tabular}
% % \end{table}

% \subsection{Reference-Link Context Enrichment}
% The previous experiments relied exclusively on the textual CVE description. However, many CVEs also reference vendor advisories, security bulletins, source-code commits, and issue trackers that contain substantially richer technical details. We therefore evaluated whether incorporating information extracted from these external references improves CWE classification performance. Up to four reference links and 4,000 characters of additional context were retrieved and appended to the prompt. This enhancement was evaluated only for the best-performing configuration identified in the previous experiments (qwen2.5-coder, Config C, temperature $=0.3$).

% Reference-link context produced the largest single improvement observed
% anywhere in this study -- a 6.0-point gain on exact-match and a 6.0-point
% gain on parent-level accuracy, exceeding the gains from taxonomy
% expansion, top-$K$ branching, model specialization, or temperature tuning
% individually. This is the best result of the entire study: 48/150 CVEs
% classified exactly right, 49/150 classified to the correct parent
% category. The cost is roughly $1.6\times$ the inference time per CVE
% (3.9s $\rightarrow$ 6.4s), driven by network I/O rather than additional
% model compute, and prediction diversity was unaffected (45/150 unique
% either way), indicating the improvement reflects genuinely
% better-informed predictions rather than a shift in guessing behavior.

% \begin{table}[th]
% \centering

% \setlength{\abovecaptionskip}{3pt}
% \setlength{\belowcaptionskip}{0pt}
% \caption{Effect of reference-link context enrichment, qwen2.5-coder (14B), T = 0.3, Config C (n = 150).}
% \label{tab:reflinks}
% \begin{tabular}{lrr}
% \toprule
% \textbf{Metric} & \textbf{   Without links} & \textbf{   With links} \\
% \midrule
% Exact-match accuracy  & 26.0\% & 32.0\% \\
% Parent-level accuracy & 26.7\% & 32.7\% \\
% Prediction diversity  & 45/150 & 45/150 \\
% Avg.\ time per CVE    & 3.9s   & 6.4s   \\
% \bottomrule
% \end{tabular}
% \end{table}

% \subsection{Qualitative Observations}
% The experiments revealed several consistent behavioral patterns across models. Expanding the taxonomy substantially benefited qwen2.5 (14B), which maintained high prediction diversity while improving classification accuracy. In contrast, llama3.2 (3B) exhibited a degenerate prediction strategy under Config B, repeatedly predicting CWE-79 for nearly every input. Although this occasionally improved exact-match accuracy, the increase resulted from repeatedly guessing a common weakness rather than accurately interpreting CVE descriptions.

% Introducing top-$K$ parent branching alleviated this failure mode by reducing hierarchical error propagation. For llama3.2, prediction diversity increased considerably, although overall accuracy decreased. For qwen2.5, the same architectural change improved both diversity and accuracy, indicating that larger models can effectively exploit the expanded search space.

% Finally, early implementations of Config C demonstrated that relying solely on the ranking order produced during the first classification stage reduced overall performance. Introducing an explicit arbitration step among all candidate weaknesses consistently improved prediction quality, suggesting that even mid-sized language models benefit from structured decision refinement rather than relying exclusively on their initial ranking.

% \subsection{Discussion}
% The experimental results demonstrate that both taxonomy design and pipeline architecture substantially influence LLM-based CWE classification. Expanding the taxonomy increased coverage of software weaknesses commonly assigned by NVD analysts, while top-$K$ parent branching reduced hierarchical error propagation by allowing multiple high-level weakness categories to remain under consideration. Together, these architectural improvements produced consistent gains for sufficiently capable language models.

% Model selection proved equally important. Although all 14B models significantly outperformed the 3B baseline, the domain-specialized qwen2.5-coder consistently achieved the highest classification accuracy, outperforming both its general-purpose counterpart and the reasoning-oriented phi4 model. This suggests that security- and code-oriented pretraining contributes more to CWE classification than general reasoning capability alone.

% The experiments also demonstrate that inference configuration influences performance. Moderate decoding temperatures produced small but consistent improvements over deterministic decoding, while incorporating information extracted from CVE reference links yielded the largest single improvement observed in the study. These results indicate that the primary limitation is often the amount of technical information available within the CVE description rather than the reasoning capability of the language model itself.

% Despite these improvements, the best description-based configuration achieved an exact-match accuracy of 32.0\%, indicating that fully automated CWE assignment remains a challenging problem. Nevertheless, the proposed hierarchy-aware prompting strategy, combined with domain-specialized language models and contextual enrichment, substantially improves classification performance over simpler approaches.

% Most importantly, these improvements directly enhance the proposed exploitability estimation framework. Since exploitability estimation is derived from the historical exploitation characteristics of the predicted CWE, increasing CWE classification accuracy improves the reliability of estimating the likelihood that newly disclosed vulnerabilities will be exploited before real-world exploitation evidence becomes available.





\section{Exploitability Experiment}

\subsection{Exploitability Ranking Using Historical CWE Exploitability}
\label{sec:exploitability-ranking}

\subsubsection{Overview}
``  
To evaluate the practical usefulness of historical weakness-level exploitability, we conducted a temporal exploitability ranking experiment. Unlike supervised exploitability prediction models, the objective of this experiment is to determine whether the historical exploitability associated with a software weakness (CWE) can directly prioritize newly disclosed vulnerabilities. This experiment also serves as the foundation for evaluating our LLM-based framework, where official CWE assignments can later be replaced with automatically predicted CWE labels.

For each newly disclosed CVE, the associated CWE is mapped to its historical exploitability probability computed from vulnerabilities published in previous years. CVEs are then ranked according to this historical exploitability score, and the resulting ranking is evaluated against vulnerabilities that were subsequently observed to be exploited in the wild.

\subsubsection{Methodology}

The experiment uses the official CVE--CWE mappings provided by the National Vulnerability Database (NVD). A vulnerability is considered exploited if it appears in at least one of the publicly available exploitation datasets used throughout this paper, namely CISA Known Exploited Vulnerabilities (KEV), VulnCheck, or GrayNoise.

To prevent temporal information leakage, historical exploitability is computed using only vulnerabilities published before the evaluation year. Specifically, for each CWE $c$, the historical exploitability probability is estimated as

\begin{equation}
P(\mathrm{Exploited}\mid c)=
\frac{E_c+\alpha P_0}
{N_c+\alpha},
\label{eq:historical_exploitability}
\end{equation}

where $E_c$ denotes the number of historically exploited CVEs mapped to CWE $c$, $N_c$ is the total number of historical CVEs associated with that CWE, $P_0$ is the overall historical exploitation rate across all historical CVEs, and $\alpha$ is a smoothing parameter.


The smoothing parameter prevents unreliable exploitability estimates for weaknesses with only a few historical observations. Without smoothing, a rarely observed CWE could receive an unrealistically high exploitability probability based on only one or two exploited vulnerabilities. The parameter $\alpha$ incorporates prior knowledge by adding pseudo-observations equivalent to the global historical exploitation rate. As the number of historical CVEs associated with a CWE increases, the influence of the smoothing term becomes negligible, and the estimated exploitability converges to the observed historical exploitability of that weakness. In this study, we performed a sensitivity analysis using $\alpha \in \{20,40,60\}$. The results showed that increasing the smoothing parameter from 20 to 40 improved the exploitability ranking performance, while further increasing it to 60 produced negligible changes. Therefore, the remaining experiments report the results obtained using $\alpha=40$.

Each CVE $v$ in the evaluation year is assigned the historical exploitability score of its mapped CWE. When multiple CWE mappings exist for a CVE, the maximum historical exploitability score is used. The ranking score assigned to each vulnerability is therefore defined as

\begin{equation}
\mathrm{RankScore}(v)=
\max_{c \in \mathrm{CWE}(v)}
P(\mathrm{Exploited}\mid c),
\label{eq:ranking_score}
\end{equation}

where $\mathrm{CWE}(v)$ denotes the set of CWE mappings associated with vulnerability $v$. CVEs are then ranked in descending order according to $\mathrm{RankScore}(v)$.


The experiment follows a rolling temporal evaluation strategy. Historical exploitability probabilities are computed using all CVEs published before the evaluation year, while testing is performed exclusively on CVEs published during that year. Specifically, five independent evaluations are performed for the period 2021--2025. For example, when evaluating CVEs published in 2023, historical exploitability probabilities are estimated using only vulnerabilities disclosed before 2023. This temporal evaluation prevents information leakage from future vulnerabilities while assessing the robustness of historical weakness-level exploitability over multiple years.


The ranked vulnerabilities are evaluated using the Area Under the Precision--Recall Curve (PR-AUC), Area Under the ROC Curve (ROC-AUC), Precision@$k$, Recall@$k$, and Normalized Discounted Cumulative Gain (NDCG@$k$), where $k \in \{50,100,250,500\}$. Because exploited vulnerabilities constitute only a small fraction of all disclosed CVEs, PR-AUC is considered the primary evaluation metric. The evaluation cutoffs assess the quality of the ranking produced by Equation~(\ref{eq:ranking_score}) under different remediation capacities. Specifically, Precision@$k$, Recall@$k$, and NDCG@$k$ evaluate the top $k$ highest-ranked vulnerabilities according to their historical exploitability scores. Smaller values of $k$ (50 and 100) measure the effectiveness of highly selective prioritization, whereas larger values (250 and 500) evaluate the ranking quality when organizations can investigate or remediate a broader set of vulnerabilities.


\subsubsection{Results}

Table~\ref{tab:cwe_ranking} summarizes the exploitability ranking performance obtained using official NVD CWE assignments across five temporal evaluation years (2021--2025). For each evaluation year, historical exploitability probabilities were computed exclusively from vulnerabilities published prior to the corresponding test year, ensuring a leakage-free assessment.

\begin{table*}[ht]
\centering
\caption{Temporal exploitability ranking using official NVD CWE assignments.}
\label{tab:cwe_ranking}
\scriptsize
\begin{tabular}{lrrrrrrrr}
\toprule
\textbf{Year} &
\textbf{CVEs} &
\textbf{  Exploited-CVEs} &
\textbf{ PR-AUC} &
\textbf{ ROC-AUC} &
\textbf{ P@100} &
\textbf{ R@100} &
\textbf{ NDCG@100} &
\textbf{ Coverage} \\
\midrule
2021 & 21,019 & 476 & 0.0520 & 0.6919 & 0.050 & 0.0105 & 0.0857 & 79.9\% \\
2022 & 25,793 & 448 & 0.0466 & 0.7016 & 0.050 & 0.0112 & 0.0964 & 84.5\% \\
2023 & 29,377 & 534 & 0.0401 & 0.7036 & 0.030 & 0.0056 & 0.0573 & 88.4\% \\
2024 & 40,484 & 677 & 0.0408 & 0.7161 & 0.000 & 0.0000 & 0.0688 & 93.2\% \\
2025 & 49,549 & 631 & 0.0352 & 0.7268 & 0.110 & 0.0174 & 0.0994 & 90.1\% \\
\bottomrule
\end{tabular}
\end{table*}

Overall, the results demonstrate that historical weakness-level exploitability provides a useful signal for prioritizing newly disclosed vulnerabilities. Across all five evaluation years, the proposed ranking consistently achieved ROC-AUC values between 0.692 and 0.727, indicating that historical exploitability effectively distinguishes vulnerabilities that are later exploited from those that are not.

The Precision--Recall performance ranged from 0.035 to 0.052 across the evaluation period. Although PR-AUC gradually decreased over time, this trend coincides with a substantial reduction in the prevalence of exploited vulnerabilities, from 2.26\% in 2021 to 1.27\% in 2025, making the ranking task increasingly challenging. Relative to the expected PR-AUC of a random ranking, which equals the prevalence of exploited vulnerabilities, the proposed approach consistently achieved approximately 2.2--2.8 times better performance, demonstrating that historical CWE exploitability contains meaningful information for identifying vulnerabilities that are eventually exploited in the wild.

The ranking quality also remained stable across different evaluation cutoffs. For example, in 2025, 11 of the top 100 ranked vulnerabilities (Precision@100 = 0.11) were subsequently observed to be exploited, despite an overall exploitation prevalence of only 1.27\%. Similarly, the NDCG@100 values ranged from 0.057 to 0.099 across the five evaluation years, indicating that exploited vulnerabilities were consistently concentrated near the top of the ranked lists.

Finally, the applicability of the approach improved over time as the availability of CWE mappings increased. The percentage of CVEs with usable CWE assignments grew from approximately 80\% in 2021 to over 90\% in 2025, demonstrating that historical weakness-level exploitability can be applied to the majority of newly disclosed vulnerabilities.

These results establish an upper bound on exploitability ranking performance when official NVD CWE assignments are available. In the next section, we replace the official CWE mappings with the CWE labels predicted by our hierarchy-aware LLM framework to evaluate whether automated CWE classification can provide comparable exploitability ranking performance immediately after CVE disclosure.

\section{Threats to Validity}

Like all empirical studies, this work has several limitations.

Our analysis assumes that CWE assignments in the National Vulnerability Database (NVD) accurately represent the underlying software weaknesses. Although RQ2 showed that few CWE mappings change after publication, some may still be incomplete, overly coarse, or incorrect. Similarly, exploitability was identified using CISA KEV, VulnCheck, and GrayNoise. While combining these sources improves coverage, exploited vulnerabilities absent from all three datasets may be incorrectly classified as non-exploited. Moreover, our LLM evaluation measures CWE classification accuracy rather than the end-to-end performance of a vulnerability prioritization system.

Methodological choices may also influence the results. Historical exploitability analyses considered only CWE-year combinations with sufficient observations, and exploitability rankings were restricted to CWEs satisfying minimum sample-size and statistical significance requirements. Although sensitivity analyses showed that the identified CWEs remain largely consistent under different thresholds, alternative thresholds or ranking weights may produce different rankings. Likewise, the LLM evaluation used specific prompting strategies, taxonomy configurations, and a limited set of open-source models; different prompts, models, or retrieval methods may yield different results.

Our empirical study relies on publicly disclosed CVEs from the NVD and publicly available exploitation datasets. Consequently, undisclosed vulnerabilities or exploits known only through proprietary threat intelligence are not represented, limiting generalizability to private vulnerability repositories or organization-specific environments. Similarly, the evaluated LLM framework targets software weaknesses defined by the CWE taxonomy and may require adaptation for other vulnerability taxonomies.

Finally, although we observed statistically significant temporal stability of historical exploitability at the CWE level, the correlations are moderate rather than strong. Historical exploitability should therefore be viewed as one informative signal rather than a complete predictor of exploitation. It is intended to complement other factors, including CVSS metrics, affected products, exploit availability, and threat intelligence, in future vulnerability prioritization frameworks.

\section{Conclusion and Future Work}

This paper investigated whether historical exploitability patterns at the software weakness level can support the early prioritization of newly disclosed vulnerabilities. We first conducted a large-scale empirical study of 192,788 CVEs and 768 CWE categories using exploitation evidence from CISA KEV, VulnCheck, and GrayNoise. Our analysis showed that CWE assignments are generally stable after publication and that the historical exploitability of many software weaknesses remains moderately stable across years. Furthermore, we identified several CWE categories that exhibit persistent associations with real-world exploitation, suggesting that weakness-level exploitability provides meaningful information beyond vulnerability severity alone.

Motivated by these findings, we proposed a hierarchy-aware LLM-based framework for automatically mapping newly disclosed CVEs to CWE categories. The framework leverages the hierarchical structure of the CWE taxonomy, taxonomy-aware prompting, and contextual information extracted from CVE reference links to improve classification performance. Our experiments demonstrated that these design choices substantially improve CWE classification accuracy, with the best configuration achieving an exact-match accuracy of 32.0\%. Although this accuracy is not yet sufficient for fully automated CWE assignment, the results demonstrate that current LLMs can provide useful assistance for vulnerability analysis and enable the use of historical CWE exploitability immediately after vulnerability disclosure.

Overall, this work establishes a practical framework that combines empirical exploitability knowledge with LLM-based vulnerability analysis to support earlier and more informed vulnerability prioritization. Rather than relying solely on CVSS scores or CVE-specific features, the proposed approach incorporates historical exploitability characteristics associated with the underlying software weakness, providing an additional source of information when exploitability evidence is unavailable.

Future work will focus on improving both CWE classification and exploitability estimation. We plan to evaluate larger and security-specialized language models, retrieval-augmented generation using additional vulnerability sources such as vendor advisories, proof-of-concept exploits, and source-code commits, and multi-label CWE prediction for vulnerabilities associated with multiple weaknesses. We also plan to integrate historical CWE exploitability with complementary signals such as CVSS metrics, EPSS scores, affected products, and threat intelligence to develop more accurate exploitability prediction models and evaluate their effectiveness on future vulnerability disclosures.




\bibliographystyle{splncs04}
\bibliography{references.bib}


\end{document}
